% !TEX root = ../main.tex
\chapter{手势识别服务实现与实验分析}


\setlength{\tabcolsep}{3pt}
HearBridge 系统中的手势识别功能主要分为两类任务：一类是面向移动端训练场景的实时单词识别，另一类是面向短视频辅助交流场景的句子视频识别。
实时单词识别用于支持用户在训练页面中观看动作示范后进行模仿练习，系统根据用户当前动作返回识别结果和置信度信息；
句子视频识别则用于支持用户选择一段手语短视频，系统从视频中识别多个词级结果，并进一步展示中文映射和语义修正结果。
本章先围绕实时单词识别说明系统闭环实现，再单独展开句子视频识别的资源选择、方法实现与实验分析。

% !TEX root = ../../main.tex
\section{实时单词识别模型与系统闭环实现}

实时单词识别是 HearBridge 系统中直接服务于移动端训练页面的识别功能。
该功能的输入来自用户在手机端实时完成的单个手势动作，输出为当前动作类别、置信度以及模型版本等信息。
与句子视频识别不同，实时单词识别不处理一段视频中的多个连续动作，也不涉及中文映射和语义修正，其主要目标是为用户模仿练习提供即时识别反馈。

% !TEX root = ../../../main.tex

\subsection{数据资源与技术选择}

实时单词识别模型主要使用系统自采样本进行训练。
管理员在系统维护过程中可通过样本采集功能为指定手势动作采集训练样本，Python 手势识别服务接收采集帧并保存为 raw dataset。
选择自采数据作为实时单词识别的数据来源，是因为该任务需要直接服务于 HearBridge 当前维护的手语资源和训练动作，自采样本更贴近系统实际训练项、采集方式和运行环境。

在技术选择上，实时单词识别采用基于关键点特征的固定词表分类方案。
识别服务使用 MediaPipe 对手部和人体姿态关键点进行提取，将原始图像帧转换为结构化时序数据，再通过特征转换流程生成模型训练所需的数据。
模型训练部分采用 TensorFlow/Keras 实现固定词表分类模型，具体使用 1D CNN 对固定长度的时序特征进行分类。

% !TEX root = ../../../main.tex

\subsection{自采样本采集链路}


\WhuAutoFigure{0.82\textwidth}{0.42\textheight}{figures/chapter06/fig6-1-raw-dataset-flow.png}{自采样本采集流程图}{fig:chapter06-raw-dataset-flow}


自采样本采集链路是实时单词识别闭环的起点。
管理员或维护人员在采集前先确定当前样本所属的手势标签，移动端随后与 Python 手势识别服务建立数据采集 WebSocket 连接，并发送采集开始消息。
服务端接收到 start\_dataset 消息后读取其中的 label 字段，如果标签为空则返回错误提示；
若标签有效，则清空当前连接中的采集缓存，将采集状态设置为 ready，并向前端返回当前标签、已采集有效帧数和目标窗口帧数等状态信息。

采集过程采用连接级缓存设计。
Python 服务在每个 WebSocket 连接中维护 raw\_frames、pending\_raw\_frames、pending\_label 和 waiting\_save\_confirm 等状态变量。
其中，raw\_frames 用于保存当前正在采集的有效帧；
pending\_raw\_frames 用于保存已经采满但尚未确认保存的样本；
pending\_label 用于记录待保存样本的标签；
waiting\_save\_confirm 用于避免样本采满后继续接收新帧导致数据混入。

移动端持续向 WebSocket 发送图像二进制数据。
服务端在接收到二进制帧后，首先判断采集是否已经启动，若尚未发送 start\_dataset 则返回错误提示；
若当前已有等待确认保存的样本，则暂不继续处理新帧。
随后，服务端根据采样时间间隔过滤过密帧，避免短时间内接收过多重复图像。

每一帧图像在服务端会先转换为 OpenCV 可处理的图像格式，再由 BGR 转为 RGB 后分别送入 MediaPipe Hands 和 MediaPipe Pose。
服务端随后提取手部关键点、手部世界坐标、手部存在状态、姿态关键点、姿态存在状态、时间戳和画面宽高等信息，整理为 raw frame。
若当前帧未检测到有效手部或必要姿态信息，服务端不会将该帧加入样本，而是向前端返回 invalid\_frame 状态。

当有效帧数量达到配置的 raw 样本窗口长度后，服务端不会立即写入磁盘，而是将当前帧序列复制到 pending\_raw\_frames，记录待保存标签，并返回 ready\_to\_save 状态。
前端若发送 save\_dataset 消息，服务端会检查当前是否存在等待保存的完整样本，并调用 raw dataset 保存工具写入样本文件；
若前端发送 discard\_dataset，则丢弃待保存样本。

该采集链路的一个重要特点是先保存 raw dataset，而不是直接保存最终训练特征。
raw dataset 中保留了后续特征转换所需的基础信息，因此当特征构造策略发生调整时，可以基于已有 raw 样本重新生成训练特征，而不必重新采集全部动作样本。

% !TEX root = ../../../main.tex

\subsection{特征数据构成}

实时单词识别模型以 MediaPipe 检测得到的结构化关键点作为训练输入。
系统在服务端调用 MediaPipe Hands 获取左右手关键点，调用 MediaPipe Pose 获取上肢姿态关键点，并将连续帧关键点转换为固定长度的时序特征。
关键点特征能够弱化背景、光照和人物外观差异对识别结果的影响，同时降低输入维度，适合移动端训练反馈场景下的实时识别任务。

本文使用左右手 21 点手部关键点，以及身体上肢的肩、肘、腕关键点作为特征来源。
手部关键点用于描述手型结构、手指弯曲状态和手部局部姿态；
上肢姿态关键点用于描述手腕、手肘相对于肩部中心的位置关系。
实时单词识别所使用的关键点选取如图~\ref{fig:chapter06-realtime-feature-166} 所示。

\WhuAutoFigure{0.82\textwidth}{0.42\textheight}{figures/chapter06/realtime-feature-source-keypoints.png}{实时单词识别特征来源关键点示意图}{fig:chapter06-realtime-feature-166}

手部坐标特征来源于 MediaPipe Hands 输出的 hand world landmarks。
该结果为每只手提供 21 个三维关键点，坐标以米为单位，原点位于手部几何中心附近。
系统对 hand world landmarks 进行尺度归一化，使不同手部大小和不同拍摄距离下的坐标数值具有更接近的尺度范围。

姿态辅助特征来源于 MediaPipe Pose 输出的 pose landmarks。
系统从中选取左右肩、左右肘和左右腕的二维归一化图像坐标，其中 $x$ 和 $y$ 坐标分别按图像宽度和高度归一化到 0 到 1 范围内。
通过上肢关键点，模型能够获得手部动作相对于身体中心的大致位置关系。

在特征组织方式上，系统将一个实时单词识别样本表示为固定长度的时序窗口。
当前模型以连续 30 帧有效图像作为一个输入样本，每帧对应 166 维特征，因此单个样本的输入形状如公式~\ref{eq:chapter06-realtime-sample-shape} 所示。

\begin{equation}
X \in \mathbb{R}^{30\times166}
\label{eq:chapter06-realtime-sample-shape}
\end{equation}

其中，$X$ 表示单个实时单词识别样本，30 表示一个样本包含 30 帧连续有效图像，166 表示每一帧的特征维度。
对于第 $t$ 帧图像，其单帧特征 $f_t$ 由左手特征、右手特征和姿态辅助特征三部分拼接得到，如公式~\ref{eq:chapter06-realtime-frame-feature} 所示。

\begin{equation}
f_t = [H^L_{78},\ H^R_{78},\ P_{10}]
\label{eq:chapter06-realtime-frame-feature}
\end{equation}

表~\ref{tab:chapter06-realtime-feature-composition} 给出了单帧 166 维特征的组成方式。
左右手特征采用相同的维度结构，均由尺度归一化后的手部坐标特征以及手指骨骼角度特征组成；
姿态辅助特征由上肢关键点相对肩中心的位置特征和肘部角度特征组成。

\begin{longtable}{L{0.15\textwidth} L{0.08\textwidth} L{0.25\textwidth} L{0.27\textwidth} L{0.15\textwidth}}
  \caption{单帧 166 维特征构成表}
  \label{tab:chapter06-realtime-feature-composition}\\
  \toprule
  特征部分 & 维度 & 数据来源 & 计算方式 & 作用 \\
  \midrule
  左手坐标特征 & 63 & MediaPipe Hands 的左手 hand world landmarks & 以手腕到中指掌指关节点的距离作为尺度因子，对 21 个三维关键点归一化后展开 & 描述左手空间形态 \\
  左手角度特征 & 15 & 左手 hand world landmarks 的骨骼连接关系 & 基于相邻骨骼向量计算夹角余弦 & 描述左手手指弯曲状态 \\
  右手坐标特征 & 63 & MediaPipe Hands 的右手 hand world landmarks & 以手腕到中指掌指关节点的距离作为尺度因子，对 21 个三维关键点归一化后展开 & 描述右手空间形态 \\
  右手角度特征 & 15 & 右手 hand world landmarks 的骨骼连接关系 & 基于相邻骨骼向量计算夹角余弦 & 描述右手手指弯曲状态 \\
  姿态位置特征 & 8 & MediaPipe Pose 的左右肩、左右肘、左右腕二维坐标 & 以肩中心为参考点进行相对化，并按肩宽进行尺度归一化 & 描述手部相对身体的位置 \\
  姿态角度特征 & 2 & MediaPipe Pose 的左右肩、肘、腕二维坐标 & 计算左右肘部夹角余弦 & 描述上肢弯曲状态 \\
  合计 & 166 & -- & 左右手特征与姿态辅助特征拼接 & 单帧输入特征 \\
  \bottomrule
\end{longtable}

对于单只手，设第 $i$ 个 hand world landmark 为 $h_i=(x_i,y_i,z_i)$，手腕关键点为 $h_0$，中指掌指关节点为 $h_9$。
系统使用手腕点到中指掌指关节点之间的欧氏距离作为手部尺度因子，如公式~\ref{eq:chapter06-hand-scale-factor} 所示。

\begin{equation}
s_h = \lVert h_9 - h_0 \rVert_2
\label{eq:chapter06-hand-scale-factor}
\end{equation}

其中，$s_h$ 表示手部尺度因子。
若尺度因子过小，系统使用极小值 $\varepsilon$ 进行保护，避免除零问题。
归一化后的手部关键点如公式~\ref{eq:chapter06-hand-landmark-normalization} 所示。

\begin{equation}
\hat{h}_i = \frac{h_i}{\max(s_h,\varepsilon)},\quad i=0,1,\ldots,20
\label{eq:chapter06-hand-landmark-normalization}
\end{equation}

其中，$\hat{h}_i$ 表示尺度归一化后的手部关键点。
归一化后的 21 个三维关键点按固定顺序展开，形成 63 维单手坐标特征。

在坐标特征基础上，系统根据 MediaPipe Hands 的手部拓扑结构构造骨骼向量，并对每根手指上的相邻骨骼向量计算夹角余弦。
对于三个连续关键点 $a$、$b$、$c$，以 $b$ 为夹角顶点。
令 $v_1=p_a-p_b$，$v_2=p_c-p_b$，则相邻骨骼向量之间的夹角余弦如公式~\ref{eq:chapter06-finger-angle-cosine} 所示。

\begin{equation}
\cos\theta = \frac{v_1 \cdot v_2}{\lVert v_1\rVert \cdot \lVert v_2\rVert + \varepsilon}
\label{eq:chapter06-finger-angle-cosine}
\end{equation}

每根手指沿骨骼链计算若干相邻骨骼夹角，五根手指共得到 15 维角度特征。
单只手由 63 维坐标特征和 15 维角度特征共同组成 78 维手部特征。

姿态辅助特征用于描述手部动作相对于上半身的位置关系。
系统从 MediaPipe Pose 的 pose landmarks 中选取左右肩、左右肘和左右腕等上肢关键点，先根据左右肩关键点计算肩中心，再将左右手腕、左右手肘转换为相对于肩中心的二维坐标，并按肩宽进行尺度归一化。
设左肩和右肩关键点分别为 $S_l$ 和 $S_r$，肩中心 $C_s$ 的计算方式如公式~\ref{eq:chapter06-shoulder-center} 所示。

\begin{equation}
C_s = \frac{S_l + S_r}{2}
\label{eq:chapter06-shoulder-center}
\end{equation}

身体尺度因子 $s_p$ 使用左右肩之间的欧氏距离表示，如公式~\ref{eq:chapter06-body-scale-factor} 所示。

\begin{equation}
s_p = \lVert S_r - S_l \rVert_2
\label{eq:chapter06-body-scale-factor}
\end{equation}

设上肢关键点集合为 $P_i\in\{W_l,W_r,E_l,E_r\}$，其中 $W_l$、$W_r$、$E_l$、$E_r$ 分别表示左腕、右腕、左肘和右肘关键点，则相对化和尺度归一化后的姿态位置特征如公式~\ref{eq:chapter06-pose-relative-normalization} 所示。

\begin{equation}
\hat{P}_i = \frac{P_i - C_s}{\max(s_p,\varepsilon)},\quad P_i\in\{W_l,W_r,E_l,E_r\}
\label{eq:chapter06-pose-relative-normalization}
\end{equation}

本文取上述 4 个上肢关键点的二维相对坐标，因此姿态位置特征共 8 维。
系统还计算左右肘部夹角，用于描述上肢弯曲状态。
以左侧肘部为例，令 $u_l=S_l-E_l$，$v_l=W_l-E_l$，则左肘夹角余弦如公式~\ref{eq:chapter06-left-elbow-angle-cosine} 所示。

\begin{equation}
\cos\theta_l = \frac{u_l \cdot v_l}{\lVert u_l\rVert \cdot \lVert v_l\rVert + \varepsilon}
\label{eq:chapter06-left-elbow-angle-cosine}
\end{equation}

右侧肘部夹角 $\cos\theta_r$ 的计算方式与左侧相同。
姿态辅助特征由 8 维相对位置特征和 2 维肘部角度特征组成。
最终，单帧 166 维特征可以表示为公式~\ref{eq:chapter06-complete-frame-feature-166}。

\begin{equation}
f_t = [L^{coord}_{63},\ L^{angle}_{15},\ R^{coord}_{63},\ R^{angle}_{15},\ P^{rel}_{8},\ P^{angle}_{2}]
\label{eq:chapter06-complete-frame-feature-166}
\end{equation}

其中，$L^{coord}_{63}$ 和 $R^{coord}_{63}$ 分别表示左、右手尺度归一化后的 63 维坐标特征，$L^{angle}_{15}$ 和 $R^{angle}_{15}$ 分别表示左、右手的 15 维手指骨骼角度特征，$P^{rel}_{8}$ 表示 8 维上肢相对位置特征，$P^{angle}_{2}$ 表示 2 维肘部角度特征。

综上，实时单词识别模型的输入特征由手部局部结构特征和身体上肢姿态辅助特征共同组成。
左手和右手的 78 维特征主要反映经过尺度归一化后的手型结构及手指弯曲状态，10 维姿态辅助特征补充手腕、手肘与肩部中心之间的相对空间关系。
通过将连续 30 帧单帧特征按时间顺序堆叠，系统最终得到形状为 $30\times166$ 的时序输入样本，为后续 1D CNN 模型进行固定词表手势分类提供数据基础。

% !TEX root = ../../../main.tex

\subsection{特征转换与训练数据生成}

在完成 raw dataset 采集后，系统需要将原始样本转换为模型可直接读取的训练数据。
自采样本保存的是连续帧中的手部关键点、姿态关键点和检测状态等原始结构化数据，而模型训练需要的是形状固定、类别明确、数值类型统一的特征样本。
因此，特征转换过程的作用就是将 raw dataset 中按类别保存的样本文件，批量转换为 $30\times166$ 的特征数组，并按照手势类别生成训练标签。
该步骤处于样本采集和模型训练之间，既要保留样本类别语义，又要将不同帧中的关键点整理为统一的数据矩阵。
实时单词识别特征转换流程如图~\ref{fig:chapter06-realtime-feature-convert-flow} 所示。

特征转换脚本首先扫描 raw dataset 根目录。
raw dataset 按手势类别分目录保存，因此脚本会依次遍历每个类别目录，并读取目录下的 sample\_XXX.npz 样本文件。
每读取一个 raw 样本，系统都会检查其帧数是否符合窗口长度要求。
当前实时单词识别模型要求每个样本包含 30 帧，如果某个样本帧数不符合要求，转换流程会跳过该样本并记录失败原因。
这种检查可以避免异常采集片段进入训练集，从而降低后续训练阶段因输入形状不一致导致失败的风险。

单个 raw 样本转换时，系统会逐帧构造 166 维特征。
每一帧先从左右手槽位中读取 hand world landmark 和手部存在状态：若某只手存在，则根据该手的 21 个三维关键点构造 78 维手型特征；
若某只手不存在，则该手对应位置填充为 0。
随后，系统根据姿态关键点构造 10 维姿态辅助特征。
每一帧的左手特征、右手特征和姿态辅助特征按固定顺序拼接，连续 30 帧再按时间顺序堆叠为一个样本。
因此，转换后的每个样本均具有 $30\times166$ 的固定形状，可以直接作为一维时序分类模型的输入。

\WhuAutoFigure{0.82\textwidth}{0.36\textheight}{figures/chapter06/realtime-feature-convert-flow.png}{实时单词识别特征转换流程图}{fig:chapter06-realtime-feature-convert-flow}

表~\ref{tab:chapter06-realtime-feature-conversion-io} 概括了特征转换阶段的输入、处理和输出。
其中，raw 样本仍然保留较完整的检测结果，feature 样本则只保留模型训练所需的数值特征。
这种分层保存方式使系统既可以回溯原始采集数据，也可以直接复用已经转换好的训练数据。

\begin{table}[H]
  \centering
  \caption{实时单词识别特征转换输入输出表}
  \label{tab:chapter06-realtime-feature-conversion-io}
  \begin{tabular}{L{0.18\textwidth} L{0.30\textwidth} L{0.42\textwidth}}
    \toprule
    环节 & 文件或数据 & 说明 \\
    \midrule
    输入目录 & \path{dataset_raw_phone_10fps/<label>/sample_XXX.npz} & 按手势标签分目录保存移动端采集得到的 raw 样本 \\
    帧数检查 & 30 帧窗口 & 只接受满足窗口长度要求的样本，异常样本记录失败原因并跳过 \\
    单帧特征 & 166 维向量 & 由左手 78 维、右手 78 维和姿态辅助 10 维特征拼接得到 \\
    输出目录 & \path{data_processed_arm_pose_10fps/<label>/sample_XXX.npy} & 保存可直接用于模型训练的定长特征样本 \\
    转换摘要 & 扫描数、转换数、跳过数和失败项 & 返回给管理端用于展示转换结果和排查异常样本 \\
    \bottomrule
  \end{tabular}
\end{table}

从工程实现角度看，特征转换并不是简单的文件格式改写。
系统在转换过程中会进行数据可用性检查、缺失关键点填充、手部尺度归一化和姿态相对化处理。
左右手不存在时使用零向量占位，可以保证样本维度稳定；
姿态点不可用或肩部尺度过小时也会采用默认值保护，避免异常数值传播到训练数据中。
完成转换后，管理端可以基于转换摘要判断本次数据集是否已经具备训练条件，并进一步触发固定词表模型训练流程。

% !TEX root = ../../../main.tex

\subsection{固定词表模型训练}

在完成特征转换后，实时单词识别模型使用生成的 $30\times166$ 时序特征进行固定词表分类训练。
训练脚本首先从特征目录中读取所有 .npy 特征样本，并根据目录名称生成对应的类别标签。
读取完成后，系统得到模型训练所需的特征矩阵 X 和标签向量 y，其中 X 的形状为“样本数 $\times$ 30 $\times$ 166”。
为了降低样本顺序对训练结果的影响，系统按照类别进行分层打乱和划分，使训练集与验证集都尽量覆盖当前固定词表中的各个手势类别。

\WhuAutoFigure{0.82\textwidth}{0.42\textheight}{figures/chapter06/realtime-cnn-structure.png}{实时单词识别 1D CNN 模型结构图}{fig:chapter06-realtime-cnn}

实时单词识别模型采用 1D CNN 结构。
由于输入数据已经是按时间排列的关键点特征序列，模型不需要从二维图像中学习空间纹理，而是需要从连续帧特征中学习动作变化模式。
因此，模型使用一维卷积沿时间维度提取局部时序特征。
模型结构如图~\ref{fig:chapter06-realtime-cnn} 所示。
输入层接收连续 30 帧、每帧 166 维的特征序列；
前两层一维卷积负责捕捉短时间范围内手型和上肢姿态的变化；
池化层和全局平均池化层用于压缩时间维度上的冗余信息；
全连接层进一步整合时序特征，并通过 Softmax 输出每个固定词表类别的概率。

表~\ref{tab:chapter06-realtime-model-train-config} 给出了实时单词识别模型的主要结构与训练参数。
该模型规模较小，能够在当前自采小词表数据上快速完成训练，并且推理阶段只需要处理关键点特征而非原始图像像素，适合用于移动端训练页面的实时反馈。

\begin{table}[H]
  \centering
  \caption{实时单词识别模型结构与训练参数表}
  \label{tab:chapter06-realtime-model-train-config}
  \begin{tabular}{L{0.20\textwidth} L{0.28\textwidth} L{0.42\textwidth}}
    \toprule
    项目 & 设置 & 作用 \\
    \midrule
    输入形状 & $30\times166$ & 表示连续 30 帧、每帧 166 维关键点特征 \\
    时序特征提取 & 两层 Conv1D & 沿时间维度提取局部动作变化模式 \\
    时间维压缩 & MaxPooling1D 与 GlobalAveragePooling1D & 降低时间维冗余并汇聚整体动作特征 \\
    分类层 & Dense + Softmax & 输出固定词表中各类别的预测概率 \\
    优化器 & Adam & 用于模型参数迭代更新 \\
    损失函数 & sparse categorical crossentropy & 适用于整数类别标签的多分类任务 \\
    训练控制 & 最多 50 轮，batch size 为 8，早停耐心值为 8 & 验证集损失长期不下降时提前停止训练并恢复较优权重 \\
    输出产物 & 模型文件、标签映射、训练曲线、混淆矩阵、评估结果和训练摘要 & 支撑后续模型登记、发布和实验分析 \\
    \bottomrule
  \end{tabular}
\end{table}

模型训练使用 Adam 优化器和稀疏多分类交叉熵损失函数，评价指标为分类准确率。
训练过程中，系统将验证集损失作为早停依据，当验证集损失在若干轮内不再下降时提前停止训练，并恢复验证集表现较好的模型权重。
训练结束后，Python 服务会保存模型文件和标签映射文件，同时生成训练曲线图、混淆矩阵图、评估结果文本和训练摘要 JSON。
这些训练产物不仅用于分析模型效果，也会被后端登记为模型版本，为管理端发布和实时识别服务重载提供依据。
因此，固定词表模型训练阶段连接了前一节的 feature 数据集和后一节的模型发布链路，是实时单词识别闭环中的核心环节。

% !TEX root = ../../../main.tex

\subsection{模型发布、重载与实时识别联动}


\WhuAutoFigure{0.82\textwidth}{0.42\textheight}{figures/chapter06/fig6-4-model-release-reload-flow.png}{模型发布、重载与实时识别联动流程图}{fig:chapter06-model-release-reload-flow}


固定词表模型训练完成后，训练产物还需要进入系统运行链路，才能真正被移动端实时识别功能使用。
HearBridge 系统通过模型版本管理、模型发布和识别服务重载机制，将训练产物接入后台管理端和实时识别服务。

模型发布过程中，管理员在后台管理端选择某个模型版本并点击发布。
管理端向 Spring Boot 服务端发送发布请求，服务端根据模型版本 ID 查询对应记录，并检查模型文件和标签映射文件是否可用。
随后，服务端调用 Python 手势识别服务提供的重载接口，使 Python 服务加载指定模型文件和标签映射。

实时识别过程中，Python 服务根据当前加载的模型输出预测类别和置信度，并将结果通过 WebSocket 返回移动端。
返回结果中除了识别标签和置信度外，还可以包含当前模型版本名称、是否使用发布模型等状态信息。


% !TEX root = ../../main.tex
\section{实时单词识别运行效果分析}

% !TEX root = ../../../main.tex

\subsection{自采数据规模与词表说明}

实时单词识别模型的训练数据来源于系统维护阶段采集的自采手势样本。
该部分数据与 HearBridge 移动端训练资源直接对应，每一个手势类别都对应系统中的一个训练动作。

\begin{longtable}{L{0.10\textwidth} L{0.18\textwidth} L{0.20\textwidth} L{0.16\textwidth} L{0.28\textwidth}}
  \caption{自采手势类别与样本数量统计表}
  \label{tab:chapter06-self-collected-samples}\\
  \toprule
  序号 & 手势标签 & 中文含义 & 样本数量 & 说明 \\
  \midrule
  1 & 待补充 & 待补充 & 待补充 & 对应系统中的一个训练动作 \\
  2 & 待补充 & 待补充 & 待补充 & 对应系统中的一个训练动作 \\
  合计 & 待补充 & 待补充 & 待补充 & 自采固定词表训练样本总数 \\
  \bottomrule
\end{longtable}

% !TEX root = ../../../main.tex

\subsection{模型训练结果}


\WhuAutoFigure{0.82\textwidth}{0.42\textheight}{figures/chapter06/fig6-5-realtime-training-curve.png}{实时单词识别训练过程损失与准确率变化曲线}{fig:chapter06-realtime-training-curve}


\begin{longtable}{L{0.22\textwidth} L{0.20\textwidth} L{0.48\textwidth}}
  \caption{实时单词识别模型评估结果}
  \label{tab:chapter06-realtime-model-evaluation}\\
  \toprule
  指标 & 结果 & 说明 \\
  \midrule
  验证集损失 & 待补充 & 由模型评估结果文件获得 \\
  验证集准确率 & 待补充 & 当前自采数据范围内的分类准确率 \\
  训练样本数 & 待补充 & 参与训练的数据量 \\
  验证样本数 & 待补充 & 参与验证的数据量 \\
  支持类别数 & 待补充 & 当前固定词表中的手势类别数量 \\
  \bottomrule
\end{longtable}


\WhuAutoFigure{0.78\textwidth}{0.48\textheight}{figures/chapter06/fig6-6-realtime-confusion-matrix.png}{实时单词识别模型混淆矩阵}{fig:chapter06-realtime-confusion-matrix}

% !TEX root = ../../../main.tex

\subsection{移动端实时识别效果}


\WhuAutoFigure{0.46\textwidth}{0.62\textheight}{figures/chapter06/mobile-realtime-result.png}{移动端实时识别效果图}{fig:chapter06-mobile-realtime-result}


模型训练完成并发布后，需要进一步验证其在移动端实时训练场景中的运行效果。
用户进入手势训练页面后，首先查看当前手语资源的动作示范，然后点击开始识别。
移动端建立与 Python 手势识别服务之间的 WebSocket 连接，并持续发送摄像头采集到的图像帧。
识别服务在接收到足够有效帧后构造特征窗口，调用当前发布模型进行推理，并将识别标签、置信度、有效帧数量和模型版本名称等信息返回移动端。

\begin{longtable}{L{0.10\textwidth} L{0.18\textwidth} L{0.18\textwidth} L{0.16\textwidth} L{0.14\textwidth} L{0.24\textwidth}}
  \caption{移动端实时识别结果示例}
  \label{tab:chapter06-mobile-realtime-examples}\\
  \toprule
  序号 & 测试动作 & 识别结果 & 置信度 & 是否正确 & 说明 \\
  \midrule
  1 & 待补充 & 待补充 & 待补充 & 待补充 & 动作识别结果示例 \\
  2 & 待补充 & 待补充 & 待补充 & 待补充 & 动作识别结果示例 \\
  \bottomrule
\end{longtable}

% !TEX root = ../../../main.tex

\subsection{当前问题与改进方向}

实时单词识别模型已经能够支撑当前固定词表下的移动端训练反馈，但该部分功能仍然存在一定局限。
首先，自采数据主要来自系统维护阶段的受控采集，采集人数、拍摄环境和动作变化范围有限，因此模型效果主要反映当前系统数据范围内的识别能力。
其次，部分手势动作可能在手型、手部位置或动作轨迹上较为接近，容易在识别过程中出现混淆。
后续优化重点应放在样本规模扩充、易混淆动作分析、采集规范优化和实时反馈体验提升等方面。


% !TEX root = ../../main.tex
\section{句子视频识别任务与数据资源选择}

实时单词识别主要用于移动端训练反馈，而句子视频识别则进一步面向短视频辅助交流场景。
该任务的目标是在有限词表条件下，从一段手语短视频中识别出多个可能连续出现的词级动作，并将识别结果整理为用户能够理解的文本内容。

% !TEX root = ../../../main.tex

\subsection{句子视频识别任务定位}

本项目中的句子视频识别并不等同于开放域连续手语翻译，而是在有限词表和受限场景下对短视频识别链路进行验证。
系统希望完成的是一条可接入 HearBridge 的工程流程：用户选择本地手语视频，识别服务输出词级候选结果，服务端组织中文映射，并进一步通过语义修正得到更自然的展示结果。

与 6.1 节中的实时单词识别相比，句子视频识别的输入形式、处理目标和输出结构都发生了变化。
实时单词识别面对的是移动端摄像头中的连续图像流，重点在于对当前训练动作进行即时反馈；
句子视频识别面对的是用户已经录制或选择的一段短视频，重点在于从完整视频中定位若干可能连续出现的词级动作，并按时间顺序组织为 Gloss 序列。
因此，本节所讨论的“句子”并不是自然语言意义上的完整句法生成结果，而是由有限词表中的多个词级识别结果构成的短语或电报式表达。
后续的中文直译和语义自然化均建立在该 Gloss 序列之上，不能替代视觉识别模块本身。

从系统职责看，句子视频识别链路被划分为视觉识别层、业务编排层和结果展示层。
Python 识别服务负责视频解码、MediaPipe 关键点检测、词级模型推理和候选词段输出；
Spring Boot 后端负责接收移动端上传请求、调用 Python 服务、保存识别记录，并在需要时调用 DeepSeek 对展示文本进行语义自然化；
HarmonyOS 移动端负责选择或录制视频，并向用户展示原始 Gloss、中文直译和自然化文本。
这种分层方式可以保证视觉识别结果被完整保留，同时允许展示层在不覆盖原始结果的前提下提高可读性。

表~\ref{tab:chapter06-sentence-task-scope} 给出了句子视频识别任务的边界与输出层次。
其中，原始识别结果主要用于评价模型与后处理算法，中文直译用于帮助普通用户理解 Gloss 含义，自然化文本则作为辅助展示结果。
后续 6.5 节中的帧级检测缓存、滑动窗口预测、词段合并和 DeepSeek 候选修正流程，分别如图~\ref{fig:chapter06-frame-cache-flow}、图~\ref{fig:chapter06-sliding-window}、图~\ref{fig:chapter06-window-predict-to-gloss} 和图~\ref{fig:chapter06-deepseek-candidate-correction} 所示。
这些流程共同构成了从上传视频到最终展示文本的完整识别链路。

\begin{longtable}{L{0.20\textwidth} L{0.28\textwidth} L{0.44\textwidth}}
  \caption{句子视频识别任务边界说明表}
  \label{tab:chapter06-sentence-task-scope}\\
  \toprule
  层次 & 输入输出 & 任务边界 \\
  \midrule
  视频输入 & 用户上传的短视频文件 & 处理离线视频，不要求摄像头实时连续返回结果 \\
  视觉识别 & 视频帧、关键点与词级模型 & 输出词段、置信度和 Top-K 候选，不直接生成最终自然语言句子 \\
  序列组织 & 按时间排列的词段结果 & 通过过滤、合并和重叠抑制形成原始 Gloss 序列 \\
  中文映射 & Gloss 标签与中文词义表 & 将 rawSequence 映射为 rawDisplayZh 和 rawTextZh \\
  展示修正 & 原始序列与候选集合 & 只对展示文本进行自然化，保留原始视觉识别结果 \\
  \bottomrule
\end{longtable}

由表~\ref{tab:chapter06-sentence-task-scope} 可知，本项目将句子视频识别定位为“有限词表词段识别 + 展示层语义组织”的工程验证任务。
这种定位一方面避免将小词表实验夸大为开放域手语翻译系统，另一方面也便于将识别结果接入 HearBridge 的移动端、后端和管理端业务流程。
在后续实验分析中，模型性能主要通过词级分类效果、拼接句子视频的序列匹配情况以及语义修正前后的展示效果进行评价。

% !TEX root = ../../../main.tex

\subsection{数据资源与技术选择}

句子视频识别任务没有直接沿用实时单词识别中的自采数据模型。
自采数据主要对应 HearBridge 系统当前维护的训练动作，更适合验证移动端实时训练反馈和模型发布闭环；
而句子视频识别需要在一段视频中识别多个连续出现的词级动作，并进一步组织为可读文本。

本项目使用 WLASL 作为句子视频识别实验的数据来源。
WLASL 是面向美国手语孤立词识别的数据集，提供了较多 Gloss 类别和对应视频样本，可用于构建词级手语识别实验\cite{li-wlasl-2020}。
需要说明的是，本文并不直接将 WLASL 等同于中文手语语料，而是利用其公开词级视频和 Gloss 标注验证“词级识别—序列组织—中文展示”的工程链路。
在系统展示层中，Gloss 标签会进一步映射为中文词义，供 HearBridge 用户理解识别结果。

为了使实验既能覆盖常见交流意图，又能保持数据规模、训练时间和错误分析的可控性，本文从 WLASL 中选取 25 个常见词汇构成小词表。
这些词汇覆盖人称、疑问、礼貌表达、时间、评价、地点、人物身份和常用动作等类型。
相比直接使用完整大词表，25 词设置更适合作为本科毕设阶段的闭环验证：它能够组合出“you want help”“please help teacher”“what you want today”等短句样例，同时又便于对每个类别的识别错误进行人工检查。
表~\ref{tab:chapter06-wlasl-small-vocab} 给出了本次句子视频识别实验采用的小词表类别及中文含义。

\begin{longtable}{L{0.10\textwidth} L{0.18\textwidth} L{0.30\textwidth} L{0.30\textwidth}}
  \caption{句子视频识别小词表类别说明表}
  \label{tab:chapter06-wlasl-small-vocab}\\
  \toprule
  序号 & Gloss 标签 & 中文含义 & 类别类型 \\
  \midrule
  1 & bad & 不好 / 坏 & 评价词 \\
  2 & deaf & 聋 / 听障 & 身份或状态相关词 \\
  3 & family & 家人 / 家庭 & 常用名词 \\
  4 & friend & 朋友 & 常用名词 \\
  5 & go & 去 & 常用动词 \\
  6 & good & 好 & 评价词 \\
  7 & help & 帮助 & 常用动词 \\
  8 & home & 家 & 常用名词 \\
  9 & language & 语言 & 常用名词 \\
  10 & learn & 学习 & 常用动词 \\
  11 & meet & 见面 & 常用动词 \\
  12 & no & 不 / 不是 & 应答词 \\
  13 & please & 请 & 礼貌表达 \\
  14 & school & 学校 & 常用名词 \\
  15 & sorry & 对不起 & 礼貌表达 \\
  16 & teacher & 老师 & 常用名词 \\
  17 & today & 今天 & 时间词 \\
  18 & tomorrow & 明天 & 时间词 \\
  19 & want & 想要 & 常用动词 \\
  20 & what & 什么 & 疑问词 \\
  21 & who & 谁 & 疑问词 \\
  22 & why & 为什么 & 疑问词 \\
  23 & work & 工作 & 常用动词 / 名词 \\
  24 & yes & 是 / 对 & 应答词 \\
  25 & you & 你 / 你们 & 人称代词 \\
  \bottomrule
\end{longtable}

由表~\ref{tab:chapter06-wlasl-small-vocab} 可知，本次小词表并不是只选择同一语义类型的动作，而是尽量保留短句组合所需的基本成分。
例如，you、who、what 和 why 可以构成主语或疑问表达；want、go、learn、meet、help 和 work 可以提供主要动作；today 与 tomorrow 用于表达时间；please 与 sorry 用于补充礼貌或情绪语义。
这种词表设计使后续拼接句子视频既能形成较自然的交流片段，又不会因为类别数量过大而掩盖模型输入、窗口划分和后处理策略本身的影响。

在技术路线上，句子视频识别部分继续采用 MediaPipe Holistic 作为关键点检测工具，并使用 TensorFlow / Keras 模型完成词级分类。
选择关键点特征而不是直接使用原始 RGB 视频，主要基于三点考虑。
第一，关键点特征能够弱化背景、服装、光照和拍摄设备差异对模型的影响。
第二，关键点序列更容易与实时单词识别部分的特征工程经验衔接，便于在同一章内比较 166 维基础特征和 232 维 plus 特征的差异。
第三，关键点表示的维度较低，适合在本地开发环境中快速完成数据构建、模型训练、滑动窗口推理和多轮实验分析。

从工程接入角度看，当前句子视频识别服务采用固定的 WLASL-mini-v2-25 主线。
Python 服务默认读取 \path{features_20f_plus} 中的特征数据和标签映射，加载 \path{models_20f_plus} 中的词级分类模型，并通过 \path{/api/sentence/recognize} 接口对外提供上传视频识别能力。
该接口返回 rawSequence、rawDisplayZh、rawTextZh、segmentTopK 和 elapsedMs 等字段。
其中，rawSequence 是视觉识别得到的原始 Gloss 序列，rawDisplayZh 与 rawTextZh 用于中文直译展示，segmentTopK 保留各词段的候选标签和置信度信息，为后续 DeepSeek 候选重排序和错误分析提供依据。

图~\ref{fig:chapter06-sentence-feature-build-flow} 展示了后续特征数据集构建流程，表~\ref{tab:chapter06-plus-feature-dim} 则给出了每帧 232 维 plus 特征的组成。
因此，本节的数据与技术选择可以概括为：以 WLASL 25 词小词表提供词级视频样本，以 MediaPipe Holistic 生成可解释的关键点输入，以 20 帧 plus 特征和词级分类模型支撑后续滑动窗口句子视频识别。

% !TEX root = ../../../main.tex

\subsection{特征方案与模型输入设计}

句子视频识别任务在特征方案上没有直接沿用实时单词识别模型中的 30 帧 166 维输入，而是采用 20 帧增强特征作为词级分类模型的输入。
这样设计的主要原因在于，句子视频识别面对的是已经录制完成的短视频，视频中可能包含多个连续词汇、动作过渡帧和不同长度的词段。
如果仍然使用较长的 30 帧窗口，单个窗口更容易覆盖多个词汇或较长过渡区域，不利于后续进行词段划分。


\WhuAutoFigure{0.82\textwidth}{0.42\textheight}{figures/chapter06/fig6-8-plus-feature-232.png}{句子视频识别 $20\times232$ plus 特征构成示意图}{fig:chapter06-plus-feature-232}


\begin{longtable}{L{0.22\textwidth} L{0.12\textwidth} L{0.34\textwidth} L{0.24\textwidth}}
  \caption{句子视频识别 plus 特征维度构成表}
  \label{tab:chapter06-plus-feature-dim}\\
  \toprule
  特征部分 & 维度 & 组成 & 作用 \\
  \midrule
  base\_166 & 166 & 左手 78 + 右手 78 + 姿态 10 & 基础手型与上肢姿态 \\
  static\_plus\_28 & 28 & 姿态点、手中心、手框、手存在状态、距离 & 单帧空间增强 \\
  dynamic\_delta\_38 & 38 & 姿态点与手部代表点相邻帧差分 & 动作变化描述 \\
  合计 & 232 & 166 + 28 + 38 & 单帧 plus 特征 \\
  \bottomrule
\end{longtable}

其中，base\_166 是句子视频识别特征中的基础部分。
需要说明的是，该部分在维度组织上与实时单词识别中的 166 维输入相近，都是由左右手手型特征和上肢姿态辅助特征组成；
但二者的关键点提取链路并不完全相同。
句子视频识别部分直接使用 MediaPipe Holistic 全身模型对视频帧进行处理，并从同一次 Holistic 结果中读取 pose\_landmarks、left\_hand\_landmarks 和 right\_hand\_landmarks，再构造后续特征。

static\_plus\_28 是句子视频识别中额外增加的静态增强特征，用于补充单帧内更丰富的空间关系。
该部分包括鼻尖、左右肩、左右肘、左右腕等上半身关键点相对于肩中心的归一化坐标，两只手的中心位置，两只手的外接框宽高，左右手是否存在的标志，以及手腕到身体中心、手腕到手部中心的距离等信息。

dynamic\_delta\_38 是用于描述动作变化的动态差分特征。
系统先从每帧中提取一个 38 维动态源向量，该向量由 9 个上半身姿态点和左右手各 5 个代表性手部点组成，每个点使用二维坐标，共 19 个点、38 维。
随后，系统计算当前帧动态源向量与上一帧动态源向量之间的差值，形成动态差分特征。


% !TEX root = ../../main.tex
\section{句子视频识别模型训练}

在完成句子视频识别的数据资源选择和特征方案设计后，还需要说明词级分类模型的训练过程。
句子视频识别并不是直接依赖实时单词识别模型，也不是在上传视频时临时训练模型，而是先使用 WLASL 小词表中的词级视频离线构建训练特征，再训练一个能够识别固定词表中单个词汇动作的时序分类模型。

% !TEX root = ../../../main.tex

\subsection{特征数据集构建与输入组织}

句子视频识别模型的训练数据来自 WLASL 小词表中的词级视频。
每个原始视频通常对应一个独立手语词汇，因此训练阶段的目标是将每个词级视频转换为一个固定形状的时序特征样本。
与实时单词识别中的自采 raw dataset 不同，句子视频识别实验直接以离线视频文件作为输入。
系统需要先从视频中提取人体姿态和双手关键点，再进一步组织为模型可以训练的特征数组。

\WhuAutoFigure{0.82\textwidth}{0.42\textheight}{figures/chapter06/sentence-feature-build-flow.png}{句子视频识别特征数据集构建流程图}{fig:chapter06-sentence-feature-build-flow}

如图~\ref{fig:chapter06-sentence-feature-build-flow} 所示，句子视频识别特征数据集的构建过程可以划分为视频读取、关键点检测、动作窗口定位、帧重采样、特征构造和数据集输出几个步骤。
系统首先根据样本清单读取每个视频的文件路径和 Gloss 标签。
其中，视频路径用于定位原始 WLASL 词级视频文件，Gloss 标签用于确定该视频对应的词汇类别。
随后，系统使用 OpenCV 读取视频帧，并将每一帧送入 MediaPipe Holistic 模型进行处理。
Holistic 检测结果中包含人体姿态关键点、左手关键点和右手关键点，为后续 plus 特征构造提供原始点位来源。

在视频帧处理阶段，系统不会直接从整段视频中等间隔抽取 20 帧。
这是因为词级视频中可能包含动作开始前的静止帧、动作结束后的停顿帧以及手部尚未进入主要动作区域的过渡帧。
如果直接对整段视频进行均匀采样，采样结果可能包含较多无效帧，从而削弱模型对核心动作片段的学习效果。
因此，系统首先根据手部检测结果确定有效动作窗口，再在动作窗口内部进行重采样。

为了定位有效动作窗口，系统会为视频中的每一帧生成 hand\_flags 标记。
当某一帧检测到左手或右手关键点时，该帧的 hand\_flags 记为 True。
当左右手均未被检测到时，该帧的 hand\_flags 记为 False。
系统统计 hand\_flags=True 的帧数，用于判断该视频中是否存在足够可靠的手部动作区域。
如果有效手部帧数量达到设定阈值，系统以第一个 hand\_flags=True 的帧作为动作区域起点，以最后一个 hand\_flags=True 的帧作为动作区域终点。
随后，系统会在动作区域前后各扩展一定数量的 padding 帧，以保留动作进入和动作收尾阶段的信息。

当有效手部帧数量不足时，系统不会直接丢弃该视频。
为了保证数据集构建过程具有一定的容错能力，系统会退化为使用视频中部窗口作为采样区域。
这种回退策略可以避免由于个别帧关键点检测失败而导致样本完全无法使用。
同时，样本来源、窗口位置和采样帧信息会被记录到 sample\_index.csv 中，便于后续进行数据质量检查和问题样本追踪。

动作窗口确定后，系统从窗口范围内重采样 20 个帧下标。
重采样采用线性采样方式，从动作窗口起点到终点均匀生成 20 个采样位置，并将采样位置转换为实际帧下标。
通过这种方式，不同长度的视频片段都可以被转换为统一长度的时序输入。
对于持续时间较短的视频，重采样可以保留其主要动作变化。
对于持续时间较长的视频，重采样可以压缩冗余帧，使模型重点学习动作的整体变化趋势。

对于每一个采样帧，系统调用 plus 特征构造方法生成单帧 232 维特征。
该特征由 base\_166、static\_plus\_28 和 dynamic\_delta\_38 三部分组成。
其中，base\_166 用于描述左右手手型结构和上肢基础姿态。
static\_plus\_28 用于补充单帧内的人体相对位置、手部中心、手部范围、手部存在状态和距离信息。
dynamic\_delta\_38 用于描述相邻帧之间姿态点和手部代表点的变化趋势。
因此，一个词级视频样本最终会被组织为 $20\times232$ 的时序特征矩阵。

完成单个视频的特征构造后，系统将该样本的特征矩阵加入输入数组 $X$。
同时，系统根据该视频对应的 Gloss 标签生成类别编号，并将其加入标签数组 $y$。
在处理全部视频样本后，系统将所有样本统一保存为特征数据集文件。
其中，$X$ 保存模型训练所需的时序特征，$y$ 保存每个样本对应的类别编号。
标签映射文件用于记录类别编号与 Gloss 标签之间的对应关系，保证模型推理结果能够反查到具体词汇。

\begin{longtable}{L{0.24\textwidth} L{0.32\textwidth} L{0.34\textwidth}}
  \caption{句子视频识别特征数据集文件说明表}
  \label{tab:chapter06-sentence-feature-files}\\
  \toprule
  文件名 & 内容 & 作用 \\
  \midrule
  X.npy & 样本数 $\times$ 20 $\times$ 232 特征数组 & 模型输入 \\
  y.npy & 类别编号数组 & 模型标签 \\
  labels.json & Gloss 标签映射 & 推理结果反查 Gloss \\
  sample\_index.csv & 样本来源、动作窗口、采样帧等信息 & 数据追踪与质量检查 \\
  config.json & 特征构建配置 & 记录构建参数 \\
  \bottomrule
\end{longtable}

如表~\ref{tab:chapter06-sentence-feature-files} 所示，特征数据集由模型输入、标签、标签映射、样本索引和配置文件共同组成。
其中，X.npy 和 y.npy 是模型训练阶段直接读取的核心文件。
labels.json 保存 Gloss 标签与类别编号的映射关系，用于训练后解释模型输出结果。
sample\_index.csv 记录每个样本对应的原始视频、动作窗口范围和实际采样帧下标。
该文件可以帮助开发者检查某个样本是否存在动作窗口定位异常、有效帧不足或采样位置不合理等问题。
config.json 则保存本次特征构建使用的参数，例如目标帧数、特征维度、动作窗口 padding 和最小有效手部帧阈值等。

通过上述处理，系统将来源不同、长度不同的视频样本统一转换为固定形状的时序特征数据。
这种输入组织方式既便于后续双向 GRU 时序分类模型进行批量训练，也为句子视频识别中的滑动窗口预测和候选词输出提供了统一的数据格式基础。
同时，sample\_index.csv 和 config.json 的保存也提高了实验的可复现性，使后续模型效果分析和特征构建参数调整更加方便。
% !TEX root = ../../../main.tex

\subsection{词级分类模型结构与训练策略}

在获得 X.npy、y.npy 和 labels.json 后，训练脚本加载全部 $20\times232$ 特征样本。
其中，X.npy 保存词级视频样本的时序特征矩阵，y.npy 保存对应的类别编号，labels.json 保存类别编号与 Gloss 标签之间的映射关系。
系统根据 labels.json 中的标签数量确定模型输出层的分类类别数。
为了保证训练集、验证集和测试集中各类别分布相对均衡，系统采用按类别分层划分策略对样本进行划分。
这种方式可以避免部分类别只出现在训练集或测试集中，从而提高模型评估结果的可靠性。

\WhuAutoFigure{0.82\textwidth}{0.42\textheight}{figures/chapter06/word-classifier-structure.png}{词级分类模型结构图}{fig:chapter06-word-classifier-structure}

如图~\ref{fig:chapter06-word-classifier-structure} 所示，词级分类模型采用轻量级双向 GRU 时序分类结构。
模型输入为 $20\times232$ 的 plus 特征序列，其中 20 表示重采样后的帧数，232 表示每一帧的特征维度。
输入序列首先经过 Masking 层，用于处理可能存在的零值填充帧。
随后，模型使用两层双向 GRU 提取时间序列中的动作变化信息。
第一层双向 GRU 设置 96 个隐单元，并返回完整时间序列，以保留每个时间步的中间表示。
第二层双向 GRU 设置 64 个隐单元，并输出定长向量，用于后续分类判断。
每层 GRU 后接 Layer Normalization，以稳定中间特征分布并改善训练过程。

双向 GRU 结构能够同时利用动作序列的前向信息和后向信息。
对于词级手语动作而言，某些类别可能在动作起始阶段较为相似，但在动作结束姿态或运动方向上存在差异。
因此，双向时序建模有助于模型综合考虑完整动作片段中的前后文变化。
与直接使用单帧特征分类相比，该结构能够更好地学习手型变化、手部运动轨迹和上肢姿态变化之间的时序关系。
同时，相较于更复杂的视频端到端模型，基于关键点特征的双向 GRU 模型参数规模较小，更适合当前小词表实验和系统原型验证。

在完成 GRU 时序特征提取后，模型通过全连接层进一步压缩特征并完成类别判别。
第一层全连接层设置 128 个神经元，并使用 ReLU 激活函数增强非线性表达能力。
随后加入 Dropout 层，以降低小样本训练场景下的过拟合风险。
第二层全连接层设置 64 个神经元，同样使用 ReLU 激活函数和 Dropout 正则化。
最后，输出层的神经元数量等于当前小词表类别数，并通过 Softmax 函数输出各 Gloss 类别的概率分布。
模型最终将概率最高的类别作为 Top-1 预测结果，同时保留 Top-3 候选结果用于后续句子识别与语义修正。

\begin{longtable}{L{0.24\textwidth} L{0.24\textwidth} L{0.40\textwidth}}
  \caption{词级分类模型训练参数表}
  \label{tab:chapter06-word-model-train-params}\\
  \toprule
  参数 & 取值 & 说明 \\
  \midrule
  输入形状 & $20\times232$ & 单个词级视频样本的时序特征 \\
  模型结构 & BiGRU + Dense & 基于关键点序列的时序分类模型 \\
  优化器 & Adam & 用于模型参数优化 \\
  学习率 & $1\times10^{-3}$ & 初始学习率 \\
  损失函数 & sparse categorical crossentropy & 多分类任务损失函数 \\
  batch\_size & 8 & 适配当前小样本训练规模 \\
  epochs & 120 & 最大训练轮数 \\
  dropout & 0.35 & 降低模型过拟合风险 \\
  评价指标 & accuracy / \path{top3_accuracy} & 评估 Top-1 与 Top-3 分类效果 \\
  \bottomrule
\end{longtable}

如表~\ref{tab:chapter06-word-model-train-params} 所示，模型训练使用 Adam 优化器和稀疏多分类交叉熵损失函数。
由于标签数组 y.npy 中保存的是整数类别编号，因此损失函数采用 sparse categorical crossentropy，而不需要将标签额外转换为 one-hot 编码。
训练过程中，accuracy 用于衡量模型 Top-1 预测是否正确，\path{top3_accuracy} 用于衡量真实类别是否出现在概率最高的前三个候选类别中。
对于后续句子视频识别任务而言，Top-3 候选结果具有重要意义。
即使某个词级片段的 Top-1 预测不完全准确，只要真实词汇仍然位于候选集合中，后续的语义修正模块仍有机会结合上下文对结果进行调整。

从整体训练策略看，词级分类模型并不直接处理原始视频像素，而是基于前一阶段构建的 $20\times232$ plus 特征进行训练。
这种方式降低了背景、光照和服装等视觉因素对模型的干扰，使模型更加关注人体姿态、手部结构和动作变化本身。
同时，固定长度输入、分层数据划分、Dropout 正则化和 Top-3 评价指标共同构成了当前句子视频识别实验的训练策略。
该模型的输出结果不仅用于单个词级视频的分类评估，也为后续连续视频中的滑动窗口预测、候选 Gloss 生成和语义修正提供基础。
% !TEX root = ../../../main.tex

\subsection{模型产物保存与评估输出}

\begin{longtable}{L{0.32\textwidth} L{0.30\textwidth} L{0.30\textwidth}}
  \caption{句子视频识别模型训练产物说明表}
  \label{tab:chapter06-sentence-model-artifacts}\\
  \toprule
  产物文件 & 内容 & 用途 \\
  \midrule
  best\_wlasl\_20f\_plus\_classifier.keras & 验证集表现最好的模型 & 推理阶段优先加载 \\
  wlasl\_20f\_plus\_classifier.keras & 训练结束时保存的最终模型 & 模型备选产物 \\
  metrics.json & 测试集指标与错误样本 & 实验结果分析 \\
  train\_config.json & 训练参数、输入形状、数据划分和训练历史 & 复现实验 \\
  confusion\_matrix.csv & 测试集混淆矩阵 & 类别混淆分析 \\
  \bottomrule
\end{longtable}

训练完成后，系统会保存最终模型和最佳模型两类模型文件。
当最佳模型文件存在时，推理链路会优先使用最佳模型；
如果最佳模型文件不存在，则退回使用最终模型文件。
除模型文件外，训练脚本还会输出 metrics.json、train\_config.json 和 confusion\_matrix.csv，为后续实验分析提供依据。


% !TEX root = ../../main.tex
\section{句子视频识别方法实现}

% !TEX root = ../../../main.tex

\subsection{视频读取与帧级检测结果缓存}

句子视频识别的输入来自移动端提交的本地视频文件。
在 HearBridge 系统中，用户可以在移动端选择一段待识别视频，并通过句子视频识别链路将视频上传到服务端。
Python 手势识别服务接收上传文件后，会先完成文件校验、临时保存和视频读取等前置处理。
这些步骤为后续关键点检测、滑动窗口预测和识别结果生成提供基础。

\WhuAutoFigure{0.82\textwidth}{0.42\textheight}{figures/chapter06/frame-cache-flow.png}{句子视频识别帧级检测缓存流程图}{fig:chapter06-frame-cache-flow}

如图~\ref{fig:chapter06-frame-cache-flow} 所示，视频读取与帧级检测结果缓存是句子视频识别流程的前置环节。
该环节的目标不是直接完成词汇分类，而是将上传视频转换为后续识别流程可以复用的帧序列和 MediaPipe 检测结果序列。
系统在这一阶段完成视频文件保存、视频元信息读取、逐帧解码、关键点检测以及检测状态统计等工作。

服务端接收到上传视频后，首先会对上传文件进行基本校验。
校验内容包括文件名后缀、文件是否为空以及文件大小是否超过限制等。
当前句子视频识别链路允许处理 mp4、mov、avi、mkv 等常见视频格式。
如果文件类型不符合要求，或者上传内容为空，系统会直接返回错误信息，避免无效输入继续进入后续识别流程。
通过前置校验，可以减少异常视频文件对服务端推理流程的影响。

校验通过后，系统会为本次识别请求创建独立的临时目录。
临时目录通常按照本次请求的 request\_id 进行区分。
随后，系统将上传视频保存为该目录下的输入文件。
这种做法可以避免不同识别请求之间的文件相互覆盖，也便于在一次识别流程结束后清理本次请求产生的临时文件。
在调试场景下，系统也可以保留临时目录，便于开发者查看上传视频、帧级中间结果和识别输出文件。

视频保存完成后，系统使用 OpenCV 打开待识别视频。
如果视频无法被正常打开，说明上传文件可能损坏、格式不受支持或内容不完整。
此时系统会终止识别流程并返回错误。
当视频能够正常打开时，系统会读取视频的基础元信息。
这些元信息包括帧宽、帧高、帧率和总帧数等。
视频元信息不会直接作为模型输入参与分类，但会被写入最终识别结果，用于辅助分析滑动窗口位置和识别结果是否合理。

完成元信息读取后，系统通过 VideoCapture 顺序读取视频中的全部帧。
每一帧都会被保存到内存中的 frames 列表中。
如果读取结束后帧列表为空，系统同样会终止识别流程。
这是因为后续滑动窗口构造需要基于完整的视频帧序列进行。
只有当系统成功获得有效帧序列后，才会进入帧级关键点检测阶段。

在获得视频帧序列后，系统使用 MediaPipe Holistic 对每一帧进行人体姿态和左右手关键点检测。
Holistic 检测结果中包含 pose\_landmarks、left\_hand\_landmarks 和 right\_hand\_landmarks。
其中，pose\_landmarks 用于描述人体姿态关键点，left\_hand\_landmarks 和 right\_hand\_landmarks 用于描述左右手关键点。
这些检测结果是后续构造 $20\times232$ plus 特征的原始点位来源。

系统不会在每一个滑动窗口中重复执行 MediaPipe 检测。
相反，系统会先对整段视频的每一帧执行一次 Holistic 检测。
每一帧的检测结果会按照帧顺序加入 results 列表。
results 列表与 frames 列表一一对应。
这样，后续任意滑动窗口只需要根据帧下标从 results 中读取缓存结果，而不需要重新对同一帧运行 MediaPipe。
这种缓存方式可以显著减少重复计算，提高句子视频识别阶段的推理效率。

在缓存 MediaPipe 检测结果的同时，系统还会记录每一帧的检测状态。
pose\_flags 用于表示当前帧是否检测到人体姿态关键点。
hand\_flags 用于表示当前帧是否至少检测到一只手。
当左手或右手任意一侧关键点存在时，该帧的 hand\_flags 即为 True。
这两个标志不直接作为分类标签使用，但能够反映输入视频的关键点检测质量。

完成整段视频的帧级检测后，系统会根据 pose\_flags 和 hand\_flags 计算检测覆盖率。
pose\_ratio 表示整段视频中成功检测到人体姿态关键点的帧比例。
any\_hand\_ratio 表示整段视频中至少检测到一只手的帧比例。
这两个指标会被写入识别结果中。
如果 pose\_ratio 或 any\_hand\_ratio 较低，说明当前视频中人体姿态或手部关键点检测不稳定。
在这种情况下，后续滑动窗口特征构造和词级预测结果也可能受到影响。

帧级检测缓存还为后续动态差分特征构造提供支持。
在滑动窗口阶段，系统会从 results 中取出窗口范围内的检测结果。
对于窗口内的每一个采样帧，系统根据缓存的 pose 和 hand 关键点构造单帧 plus 特征。
同时，系统会在窗口内部维护上一帧的动态源向量，用于计算相邻帧之间的 dynamic\_delta 特征。
窗口第一帧没有上一帧可供比较，因此其动态差分可以视为从零开始。
这种处理方式与训练阶段对独立词级样本构造 20 帧特征的方式保持一致。

通过上述设计，系统将视频读取、帧级关键点检测和滑动窗口特征构造进行了解耦。
视频读取阶段负责获得帧序列和元信息。
帧级检测阶段负责一次性缓存每一帧的 MediaPipe 结果。
滑动窗口阶段则复用这些缓存结果构造多个候选窗口的模型输入。
这种流程既减少了重复检测开销，也使后续识别结果能够结合视频元信息、检测覆盖率和窗口位置进行分析。
最终，帧序列、检测缓存、pose\_ratio、any\_hand\_ratio 等信息会共同支撑后续的窗口预测、词段合并和句子级识别结果生成。
% !TEX root = ../../../main.tex

\subsection{滑动窗口特征构造与推理输入生成}


\WhuAutoFigure{0.78\textwidth}{0.42\textheight}{figures/chapter06/sliding-window.png}{滑动窗口构造示意图}{fig:chapter06-sliding-window}


在完成视频读取和逐帧 MediaPipe 结果缓存后，系统需要将整段句子视频转换为词级分类模型可以接收的批量输入。
由于输入视频中可能包含多个连续词汇，系统不会将整段视频压缩为一个整体样本，而是在完整帧序列上构造多个固定长度的候选窗口。
每个候选窗口对应一段可能包含单个手语词汇主体动作的局部片段，并被转换为形状为 $20\times232$ 的时序特征矩阵。

滑动窗口由 build\_windows 函数统一生成。
设整段视频的帧级检测结果数量为 total\_frames，窗口长度为 window\_size，步长为 stride。
默认 window\_size 为 20，stride 为 2。
若 total\_frames 小于或等于 window\_size，系统只构造一个从第 0 帧开始的窗口；
若 total\_frames 大于 window\_size，系统从第 0 帧开始，每隔 stride 帧生成一个窗口起点，直到无法继续形成完整窗口。

\begin{equation}
window\_size = 20,\quad stride = 2
\label{eq:chapter06-window-size-stride}
\end{equation}

\begin{equation}
starts = \{0,\ stride,\ 2\times stride,\ \cdots\}
\label{eq:chapter06-window-starts}
\end{equation}

若最后一个 start 不等于 $total\_frames - window\_size$，则补充 $last\_start = total\_frames - window\_size$。

单个窗口的特征由 build\_window\_feature 函数构造。
函数从指定 start 开始，连续读取 window\_size 个帧级 MediaPipe 结果，并逐帧调用 plus 特征提取函数。
每一帧会被转换为 232 维 plus 特征；
连续 20 帧堆叠后，形成一个 $20\times232$ 的候选窗口输入。

动态差分特征在每个滑动窗口内部重新计算。
构造单个窗口特征时，系统会先将 previous\_dynamic\_source 初始化为空；
窗口内第一帧由于没有上一帧作为差分参照，其 dynamic\_delta 被置为全 0；
从第二帧开始，系统才根据当前帧和前一帧的 dynamic\_source 计算相邻帧差分。

\begin{equation}
X \in \mathbb{R}^{window\_count\times20\times232}
\label{eq:chapter06-batch-infer-shape}
\end{equation}

% !TEX root = ../../../main.tex

\subsection{词级分类预测与 Gloss 序列生成}

在生成批量推理输入张量 $X$ 后，
系统将其整体输入已经训练完成的词级分类模型，
得到每个滑动窗口在全部 Gloss 类别上的预测概率分布。

由于 $X$ 中的每个样本都对应一个候选时间窗口，
因此模型输出的每一行概率向量都可以通过 \texttt{window\_meta} 映射回原视频中的具体时间区间，
包括该窗口的起始帧 \texttt{start\_frame}、结束帧 \texttt{end\_frame} 和中心帧 \texttt{center\_frame}。

从工程实现上看，
句子视频识别中的 Gloss 序列生成并不是对所有窗口预测结果直接按时间顺序拼接，
而是要依次经过逐窗口预测组织、原始词段合并、词段过滤和词段级非极大值抑制等步骤，
最终再从保留下来的稳定词段中提取标签序列。

\WhuAutoFigure{0.82\textwidth}{0.42\textheight}{figures/chapter06/window-predict-to-gloss.png}{逐窗口预测到 Gloss 序列生成流程图}{fig:chapter06-window-predict-to-gloss}

如图~\ref{fig:chapter06-window-predict-to-gloss} 所示，
系统首先根据模型输出生成逐窗口预测表 \texttt{dense\_rows}，
然后基于有效窗口构造原始词段 \texttt{raw\_segments}，
再经过词段过滤和词段级 NMS 得到最终词段集合 \texttt{final\_segments}，
最后按时间顺序提取各词段的标签，
形成句子视频识别的原始 Gloss 序列 \texttt{rawSequence}。

逐窗口预测结果由 \texttt{make\_dense\_rows} 函数统一组织。
对于第 $i$ 个候选窗口，
系统首先对类别概率按从高到低排序，
得到 Top-1、Top-2 和 Top-3 候选类别及其对应概率。

其中，
Top-1 类别表示该窗口当前最可能对应的 Gloss，
Top-2 和 Top-3 则主要用于保留候选信息，
便于后续调试分析与接口结果展示。

为了避免将低置信度窗口或类别边界不明显的窗口直接纳入后续词段构造，
系统会对每个窗口计算 Top-1 概率 \texttt{top1\_prob}，
以及 Top-1 与 Top-2 概率之间的差值 \texttt{margin}。

只有当窗口同时满足置信度阈值和类别间隔阈值时，
该窗口才会被标记为有效窗口，
即 \texttt{accepted=1}。

\begin{equation}
accepted=
\bigl(top1\_prob \geq confidence\_threshold\bigr)
\land
\bigl(margin \geq margin\_threshold\bigr)
\label{eq:chapter06-window-accepted}
\end{equation}

经过上述处理后，
系统得到逐窗口预测表 \texttt{dense\_rows}。
其中每一行对应一个滑动窗口，
记录了窗口的时间范围、中心位置、Top-1/Top-2/Top-3 预测结果以及是否被接受等信息。

需要特别说明的是，
后续词段构造主要依赖被接受窗口的 Top-1 标签，
而不是直接对 Top-3 结果做投票生成最终序列。

系统随后调用 \texttt{build\_raw\_segments}，
根据 \texttt{dense\_rows} 中所有 \texttt{accepted=1} 的窗口构造原始词段 \texttt{raw\_segments}。

具体而言，
系统按照时间顺序遍历有效窗口。
如果当前窗口与正在构造的词段具有相同的 Top-1 标签，
并且该窗口起始帧与当前词段结束帧之间的距离不超过 \texttt{same\_label\_merge\_gap}，
则将二者合并为同一个原始词段。

在合并过程中，
系统会同步更新词段的起止帧、中心帧、窗口数量、置信度累积和最大置信度等统计信息。

对于每个原始词段，
系统还会进一步计算平均置信度 \texttt{avg\_confidence} 和最大置信度 \texttt{max\_confidence}。
其中，
平均置信度反映该词段在多个连续窗口上的整体稳定程度，
最大置信度反映该词段内部最强响应窗口的概率峰值。

在得到原始词段 \texttt{raw\_segments} 后，
系统调用 \texttt{filter\_segments} 对其进行进一步筛选，
得到过滤后的词段集合 \texttt{filtered\_segments}。

该阶段的过滤规则主要包括三类。

第一，
标签为 \texttt{blank} 的词段不会进入最终 Gloss 序列。

第二，
窗口数量小于 \texttt{min\_segment\_windows} 的词段会被丢弃，
因为这类词段持续时间过短，
通常更可能是边界抖动或局部误报。

第三，
如果某个词段同时满足平均置信度较低且最大置信度也较低，
则认为该片段整体可信度不足，
同样会被过滤掉。

经过上述过滤后，
系统虽然已经去除了明显不可靠的词段，
但在短句视频中，
仍然可能保留多个时间位置接近、但标签不同的候选词段。

这类现象的根本原因在于滑动窗口之间存在较大重叠，
边界区域可能同时激活多个不同类别。
如果不进一步抑制这些竞争片段，
最终得到的 Gloss 序列中就可能出现插入误报或局部重复。

为此，
系统对 \texttt{filtered\_segments} 执行词段级非极大值抑制，
得到最终词段集合 \texttt{final\_segments}。

在当前实现中，
该过程由 \texttt{nms\_segments} 完成。
其目的不是在逐窗口层面删除预测，
而是在词段层面保留更稳定、更可信的候选结果。

NMS 的第一步是再次合并时间上相邻或重叠的同标签词段，
用于消除同一个词被拆成多个连续片段的情况。

NMS 的第二步是为每个候选词段计算稳定性分数，
并按该分数从高到低排序。
当前实现采用如下评分方式：

\begin{equation}
score = avg\_confidence \times \log(1 + window\_count)
\label{eq:chapter06-segment-score}
\end{equation}

该评分方式不仅考虑平均置信度，
还考虑该词段由多少个窗口共同支持。
因此，
与仅依据最大置信度排序相比，
该评分更有利于保留持续时间更长、概率更稳定的真实词段。

NMS 的第三步是判断不同词段之间是否存在竞争关系。
当前实现主要从两个角度进行判断。

其一是中心帧距离。
如果两个词段的中心帧距离不超过抑制半径 \texttt{nms\_suppress\_radius}，
则认为它们可能对应同一段局部动作区域。

其二是时间区间重叠程度。
如果两个词段在时间上高度重叠，
或者其中一个词段的大部分区间被另一个词段覆盖，
系统也会将它们视为竞争候选。

对于存在竞争关系的候选词段，
系统只保留稳定性分数更高的那个，
从而抑制边界误报和局部插入误报。

经过上述处理后，
系统得到按时间顺序排列的最终词段集合 \texttt{final\_segments}。
每个最终词段都包含 \texttt{label}、\texttt{start\_frame}、\texttt{end\_frame}、\texttt{window\_count}、\texttt{avg\_confidence} 和 \texttt{max\_confidence} 等信息。

最终的 Gloss 序列正是由 \texttt{final\_segments} 生成的。
设最终词段序列为

\begin{equation}
final\_segments = \{s_1, s_2, \dots, s_K\}
\label{eq:chapter06-final-segments}
\end{equation}

其中各词段已按起始帧升序排列，
则系统输出的原始 Gloss 序列 \texttt{rawSequence} 定义为

\begin{equation}
rawSequence = [\,label(s_1),\ label(s_2),\ \dots,\ label(s_K)\,]
\label{eq:chapter06-raw-sequence}
\end{equation}

也就是说，
Gloss 序列并不是直接由逐窗口预测表逐行拼接得到，
而是经过“窗口判定—原始词段合并—词段过滤—词段级 NMS”之后，
从最终保留下来的词段标签中按时间顺序依次提取而来。

在服务返回阶段，
Python 侧会进一步基于 \texttt{rawSequence} 生成中文显示序列 \texttt{rawDisplayZh} 和中文文本结果 \texttt{rawTextZh}，
但其核心识别输出仍然是上述原始 Gloss 序列。

因此，
如图~\ref{fig:chapter06-window-predict-to-gloss} 所示的整条处理链路，
本质上回答了句子视频识别中“Gloss 序列是如何生成的”这一关键问题，
即系统先从逐窗口概率结果构造稳定的词段，
再由最终词段映射为句子级 Gloss 序列输出。
% !TEX root = ../../../main.tex

\subsection{候选序列生成与 DeepSeek 语义修正}

在得到原始 Gloss 序列 \texttt{rawSequence} 后，
系统还需要将词级识别结果转换为用户能够理解的自然语言表达。

由于句子视频识别结果来自滑动窗口分类和词段级后处理，
\texttt{rawSequence} 中仍然可能存在边界插入词、局部词汇误判或表达不够自然等问题。

因此，
本系统在直接中文映射的基础上，
引入 DeepSeek 语义修正模块，
用于对后端生成的候选 Gloss 序列进行排序选择，
并将所选序列转换为忠实、自然的英文和中文表达。

需要说明的是，
DeepSeek 并不直接参与视频识别，
也不会重新分析原始图像帧。

视觉识别仍由 Python 端完成，
包括视频帧读取、MediaPipe 关键点检测、词级分类模型推理和词段后处理。

DeepSeek 接收的是后端整理后的文本序列、词段候选信息和程序生成的有限候选序列。

也就是说，
DeepSeek 的作用不是自由改写识别结果，
而是在受到候选序列约束的前提下，
选择更合理的候选序列并完成忠实自然化表达。

\WhuAutoFigure{0.82\textwidth}{0.42\textheight}{figures/chapter06/deepseek-candidate-correction-flow.png}{DeepSeek 候选序列语义修正流程图}{fig:chapter06-deepseek-candidate-correction}

如图~\ref{fig:chapter06-deepseek-candidate-correction} 所示，
语义修正流程可以分为四个阶段。

第一，
Python 识别服务返回原始 Gloss 序列、中文逐词映射结果和每个词段的 Top-K 候选信息。

第二，
后端根据原始词段和 Top-K 候选生成一组合法候选序列。

第三，
后端将原始结果、词段证据和候选序列封装为 DeepSeek 请求。

第四，
DeepSeek 从候选集合中选择一个候选编号，
后端再根据该候选编号生成最终修正结果。

表~\ref{tab:chapter06-deepseek-input-fields} 给出了 DeepSeek 语义修正请求中的主要输入字段、数据来源及其作用。

\begin{longtable}{L{0.20\textwidth} L{0.24\textwidth} L{0.28\textwidth} L{0.20\textwidth}}
  \caption{DeepSeek 语义修正输入信息说明表}
  \label{tab:chapter06-deepseek-input-fields}\\
  \toprule
  输入字段 & 数据来源 & 主要内容 & 作用 \\
  \midrule
  rawSequence & Python 句子视频识别结果 & 原始 Gloss 序列 & 表示视觉识别模块输出的词级识别结果 \\
  rawTextZh & 后端中文映射结果 & 逐词映射后的中文文本 & 作为直接展示结果和自然化翻译参照 \\
  segments & Python 返回的 segmentTopK 经后端整理得到 & 词段编号、原始标签、中文标签、置信度和 Top-K 候选 & 为候选选择提供视觉识别证据 \\
  candidates & 后端候选生成模块 & 由保留、删除和替换操作组合生成的合法候选序列 & 约束 DeepSeek 只能从程序候选中选择，避免自由改写 \\
  \bottomrule
\end{longtable}

DeepSeek 请求中的 \texttt{segments} 字段来源于 Python 侧返回的 \texttt{segmentTopK}。

每个词段不仅包含原始识别标签，
还包含该词段范围内的 Top-K 候选标签及其概率统计。

这些候选信息为后端生成替换候选提供依据，
也为 DeepSeek 判断某个替换是否具有视觉证据提供参考。

在字段结构上，
语义修正请求可以抽象为代码清单~\ref{lst:chapter06-deepseek-input-example} 所示形式。

\begin{lstlisting}[style=whuCode,caption={DeepSeek 语义修正输入结构示例},label={lst:chapter06-deepseek-input-example}]
{
  "rawSequence": ["you", "want", "learn"],
  "rawTextZh": "你 想要 学习",
  "segments": [
    {
      "segmentIndex": 1,
      "rawLabel": "you",
      "rawLabelZh": "你",
      "avgConfidence": 0.86,
      "maxConfidence": 0.94,
      "topK": [
        {
          "label": "you",
          "labelZh": "你",
          "avgProb": 0.86
        },
        {
          "label": "who",
          "labelZh": "谁",
          "avgProb": 0.21
        }
      ]
    }
  ],
  "candidates": [
    {
      "candidateId": "C000",
      "sequence": ["you", "want", "learn"],
      "edits": [],
      "editCount": 0,
      "replaceCount": 0,
      "removeCount": 0,
      "source": "raw"
    },
    {
      "candidateId": "C001",
      "sequence": ["you", "learn"],
      "edits": [
        {
          "type": "remove",
          "segmentIndex": 2,
          "from": "want",
          "to": null
        }
      ],
      "editCount": 1,
      "replaceCount": 0,
      "removeCount": 1,
      "source": "remove"
    }
  ]
}
\end{lstlisting}

后端首先基于原始词段和每个词段的 Top-K 候选构造候选 Gloss 序列。

对于每个词段，
程序会生成三类操作选项。

第一类是保留原标签，
即该词段继续使用视觉识别模块给出的 \texttt{rawLabel}。

第二类是删除该词段，
用于处理明显的边界插入词、重复词或过渡噪声。

第三类是替换该词段，
即使用同一词段 Top-K 候选中的其他标签替代原标签。

随后，
系统对各词段的操作选项进行组合，
形成候选序列集合 \texttt{candidates}。

为了避免候选空间过大，
并防止语义修正幅度过强，
后端对候选生成过程设置了约束。

当前实现中，
总编辑次数最多为 2 次，
替换次数最多为 1 次，
候选序列数量最多为 81 个。

如果候选序列为空，
或者原始序列本身为空，
系统会直接进入 fallback 分支，
不再调用 DeepSeek。

DeepSeek 请求由后端统一构造。

系统提示词将 DeepSeek 约束为手语 Gloss 候选排序器和忠实自然化翻译器，
要求其只能从 \texttt{candidates} 中选择一个 \texttt{candidateId}。

DeepSeek 不允许新增词段，
不允许重排词段，
也不允许使用候选集合之外的 Gloss 词。

同时，
DeepSeek 返回的 \texttt{correctedGlossSequence} 必须与所选候选序列完全一致。

表~\ref{tab:chapter06-deepseek-output-fields} 总结了 DeepSeek 返回结果中的主要字段含义及后端处理方式。

\begin{longtable}{L{0.24\textwidth} L{0.32\textwidth} L{0.34\textwidth}}
  \caption{DeepSeek 语义修正返回字段说明表}
  \label{tab:chapter06-deepseek-output-fields}\\
  \toprule
  返回字段 & 含义 & 系统处理方式 \\
  \midrule
  selectedCandidateId & DeepSeek 从候选集合中选择的候选编号 & 后端校验该编号是否存在于 candidates 中，并以该编号对应的程序候选序列为准 \\
  correctedGlossSequence & DeepSeek 写出的修正后 Gloss 序列 & 理论上应与 selectedCandidateId 对应候选一致；若不一致，后端仍信任 selectedCandidateId 对应候选 \\
  englishSentence & 忠实英文句子或片段 & 用于中间语义解释和调试展示 \\
  chineseTranslation & 忠实自然中文表达 & 作为 correctedTextZh 的主要来源，用于前端展示 \\
  translationConfidence & 翻译置信等级 & 用于提示用户当前自然化结果的可靠程度 \\
  isIncomplete & 语义是否不完整 & 若为 true，说明结果更适合作为片段展示 \\
  translationNote & 不确定性或缺失信息说明 & 用于解释低置信度或语义不完整情况 \\
  reason & 选择候选的简要原因 & 用于调试和分析候选选择依据 \\
  \bottomrule
\end{longtable}

一次正常的 DeepSeek 语义修正返回结果可以抽象为代码清单~\ref{lst:chapter06-deepseek-output-example} 所示形式。

\begin{lstlisting}[style=whuCode,caption={DeepSeek 语义修正返回结构示例},label={lst:chapter06-deepseek-output-example}]
{
  "selectedCandidateId": "C000",
  "correctedGlossSequence": ["you", "want", "learn"],
  "englishSentence": "You want to learn.",
  "chineseTranslation": "你想学习。",
  "translationConfidence": "high",
  "isIncomplete": false,
  "translationNote": "The gloss relation is clear.",
  "reason": "The raw candidate is coherent and supported by segment evidence."
}
\end{lstlisting}

后端接收到 DeepSeek 返回内容后，
不会无条件采用模型写出的 \texttt{correctedGlossSequence}。

系统首先解析 DeepSeek 返回的 JSON，
并读取其中的 \texttt{selectedCandidateId}。

随后，
后端在程序生成的候选集合中查找该编号对应的候选序列。

如果该编号存在，
最终修正后的 Gloss 序列以该候选序列为准。

如果 DeepSeek 写出的 \texttt{correctedGlossSequence} 与所选候选序列不一致，
后端也不会直接信任模型写出的序列，
而是继续使用 \texttt{selectedCandidateId} 对应的程序候选。

这一设计保证了语义修正结果始终被限制在后端生成的候选空间中，
避免大语言模型生成不受视觉证据支持的新词、
新增词段或重排词序。

在候选选择完成后，
后端会根据所选候选生成 \texttt{correctedGlossSequence}，
并结合 DeepSeek 返回的中文表达生成 \texttt{correctedTextZh}。

同时，
系统还会记录 \texttt{selectedSegments} 和 \texttt{removedSegments} 等辅助信息，
用于说明哪些词段被保留、替换或删除。

例如，
当某个候选序列表示删除一个边界插入噪声词时，
DeepSeek 只需要选择对应的候选编号，
后端即可根据该候选的编辑记录生成删除词段说明。

如果 DeepSeek 调用失败、
返回内容无法解析、
缺少 \texttt{selectedCandidateId}，
或返回的候选编号不存在于候选集合中，
系统会进入 fallback 分支。

fallback 分支会保留原始 \texttt{rawSequence} 和 \texttt{rawTextZh}，
并将 \texttt{selectedCandidateId} 设置为 \texttt{C000}。

此时，
系统不会标记语义修正已生效，
而是将结果视为原始识别结果的直接展示。

通过上述机制，
本系统将大语言模型限定在候选排序和忠实自然化表达范围内。

这种设计既利用了 DeepSeek 对短句语义合理性的判断能力，
又避免了大语言模型自由生成导致的识别结果失真。

因此，
候选序列生成与 DeepSeek 语义修正模块并不是替代视觉识别模块，
而是在视觉识别结果基础上，
对有限候选进行保守选择和自然语言表达转换，
从而提升句子视频识别结果的可读性与用户可理解性。

% !TEX root = ../../main.tex
\section{句子视频识别实验结果分析}

本节用于对句子视频识别模型和完整识别链路的运行效果进行分析。
具体数值、表格和图像需要根据后续整理到的 \path{metrics.json}、\path{confusion_matrix.csv}、句子视频推理结果、DeepSeek 修正前后结果等实验产物填写。
当前先保留论文结构和插入位置。

% !TEX root = ../../../main.tex

\subsection{词级分类模型测试结果}

本节对句子视频识别流程中使用的词级分类模型进行测试结果分析。

该模型使用 WLASL-mini-v2-25 数据集构建，
覆盖 25 个常用 Gloss 类别。

模型输入为 20 帧时序窗口，
每帧采用 232 维 plus 特征，
因此单个样本的输入形状为 $20\times232$。

本次实验共包含 195 个样本，
其中训练集 125 个样本，
验证集 25 个样本，
测试集 45 个样本。

表~\ref{tab:chapter06-word-model-test-result} 给出了词级分类模型在测试集上的主要指标。

\begin{longtable}{L{0.24\textwidth} L{0.22\textwidth} L{0.44\textwidth}}
  \caption{词级分类模型测试结果表}
  \label{tab:chapter06-word-model-test-result}\\
  \toprule
  指标 & 结果 & 说明 \\
  \midrule
  test\_loss & 1.1558 & 测试集损失 \\
  test\_accuracy & 0.7111 & 测试集 Top-1 分类准确率 \\
  test\_top3\_accuracy & 0.9333 & 测试集中真实标签进入 Top-3 的比例 \\
  train\_count & 125 & 训练集样本数量 \\
  val\_count & 25 & 验证集样本数量 \\
  test\_count & 45 & 测试集样本数量 \\
  class\_count & 25 & 词表类别数量 \\
  \bottomrule
\end{longtable}

从表~\ref{tab:chapter06-word-model-test-result} 可以看出，
模型在测试集上的 Top-1 准确率为 71.11\%，
Top-3 准确率为 93.33\%。

这说明模型在部分样本上仍存在首位类别误判，
但多数情况下能够将真实类别保留在前三候选之内。

对于本系统而言，
这一结果具有重要意义。

句子视频识别并不是仅依赖单个窗口的 Top-1 结果直接输出最终句子，
而是会在后续流程中结合 Top-K 候选、词段置信度和候选序列约束进行进一步处理。

因此，
较高的 Top-3 准确率表明模型能够为候选序列生成与语义修正模块提供较有价值的视觉候选基础。

图~\ref{fig:chapter06-word-model-confusion-matrix} 展示了词级分类模型在测试集上的混淆矩阵。

\WhuAutoFigure{0.78\textwidth}{0.48\textheight}{figures/chapter06/word-model-confusion-matrix.png}{词级分类模型混淆矩阵}{fig:chapter06-word-model-confusion-matrix}

结合图~\ref{fig:chapter06-word-model-confusion-matrix} 和错误样本可以看出，
模型的错误主要集中在少数相似动作或局部姿态容易混淆的类别之间。

例如，
\texttt{you} 与 \texttt{yes} 之间存在双向混淆，
\texttt{want} 容易被预测为 \texttt{today} 或 \texttt{help}，
\texttt{what} 容易被预测为 \texttt{today} 或 \texttt{language}。

此外，
\texttt{good}、\texttt{who}、\texttt{why} 等类别也出现了不同程度的误判。

这些错误说明，
在小样本词级分类场景中，
仅依靠单一 Top-1 输出容易受到局部手型、运动轨迹和动作边界的影响。

表~\ref{tab:chapter06-word-model-error-analysis} 给出了测试集中的错误样本分析结果。

\begin{longtable}{L{0.14\textwidth} L{0.16\textwidth} L{0.16\textwidth} L{0.34\textwidth} L{0.20\textwidth}}
  \caption{词级分类模型错误样本分析表}
  \label{tab:chapter06-word-model-error-analysis}\\
  \toprule
  样本编号 & 真实 Gloss & 预测 Gloss & Top-3 候选 & 错误说明 \\
  \midrule
  4 & \texttt{you} & \texttt{yes} & \texttt{yes}(0.6075), \texttt{you}(0.3631), \texttt{language}(0.0185) & 真实标签位于 Top-2，属于首位候选混淆 \\
  123 & \texttt{yes} & \texttt{you} & \texttt{you}(0.9415), \texttt{yes}(0.0448), \texttt{work}(0.0079) & \texttt{yes} 与 \texttt{you} 混淆明显 \\
  10 & \texttt{want} & \texttt{today} & \texttt{today}(0.9607), \texttt{want}(0.0328), \texttt{good}(0.0022) & 真实标签位于 Top-2，但 Top-1 置信度较高 \\
  1 & \texttt{you} & \texttt{yes} & \texttt{yes}(0.7524), \texttt{you}(0.1686), \texttt{deaf}(0.0369) & \texttt{you} 被误判为 \texttt{yes} \\
  133 & \texttt{what} & \texttt{today} & \texttt{today}(0.8864), \texttt{want}(0.1120), \texttt{family}(0.0003) & 真实标签未进入 Top-3 \\
  147 & \texttt{who} & \texttt{deaf} & \texttt{deaf}(0.9226), \texttt{who}(0.0220), \texttt{no}(0.0215) & 真实标签位于 Top-2，但概率较低 \\
  106 & \texttt{good} & \texttt{today} & \texttt{today}(0.8022), \texttt{want}(0.0816), \texttt{friend}(0.0457) & 真实标签未进入 Top-3 \\
  137 & \texttt{what} & \texttt{language} & \texttt{language}(0.9835), \texttt{what}(0.0097), \texttt{go}(0.0031) & 真实标签位于 Top-2，但概率极低 \\
  124 & \texttt{yes} & \texttt{you} & \texttt{you}(0.6760), \texttt{no}(0.1669), \texttt{yes}(0.0732) & 真实标签位于 Top-3，存在候选可恢复空间 \\
  149 & \texttt{why} & \texttt{home} & \texttt{home}(0.5075), \texttt{why}(0.4789), \texttt{deaf}(0.0074) & Top-1 与真实标签概率接近，边界较不确定 \\
  103 & \texttt{good} & \texttt{who} & \texttt{who}(0.7032), \texttt{good}(0.1184), \texttt{home}(0.0492) & 真实标签位于 Top-2 \\
  13 & \texttt{want} & \texttt{help} & \texttt{help}(0.7228), \texttt{learn}(0.1340), \texttt{school}(0.1228) & 真实标签未进入 Top-3 \\
  75 & \texttt{school} & \texttt{work} & \texttt{work}(0.5368), \texttt{help}(0.3173), \texttt{school}(0.1364) & 真实标签位于 Top-3，存在候选可恢复空间 \\
  \bottomrule
\end{longtable}

由表~\ref{tab:chapter06-word-model-error-analysis} 可知，
测试集中共有 13 个 Top-1 错误样本。

其中，
部分错误样本的真实标签仍然位于 Top-2 或 Top-3 候选中，
例如样本 4、样本 10、样本 124 和样本 75。

这类样本虽然 Top-1 判断错误，
但模型仍然在候选分布中保留了正确类别，
因此后续可以通过词段 Top-K 候选和候选序列修正机制进行一定程度的补偿。

与此同时，
也有部分样本的真实标签未进入 Top-3，
例如 \texttt{what} 被预测为 \texttt{today}、
\texttt{good} 被预测为 \texttt{today}、
\texttt{want} 被预测为 \texttt{help} 等情况。

这类错误说明模型在部分词汇之间仍存在较强混淆，
后续若继续提升识别性能，
需要从扩充样本数量、增强动作边界标注质量和优化特征表示等方面进行改进。

总体来看，
该词级分类模型在 25 类小词表测试集上取得了 71.11\% 的 Top-1 准确率和 93.33\% 的 Top-3 准确率。

虽然 Top-1 准确率仍有提升空间，
但较高的 Top-3 准确率表明模型能够为大多数测试样本提供有效候选集合。

因此，
在系统设计中，
后续句子视频识别流程并未简单依赖单窗口 Top-1 分类结果，
而是进一步结合滑动窗口词段合并、Top-K 候选保留、候选序列生成和 DeepSeek 受约束语义修正，
以提升最终句子级识别结果的稳定性和可读性。
% !TEX root = ../../../main.tex

\subsection{拼接句子视频整体识别结果}

在词级分类模型测试的基础上，
本文进一步对拼接句子视频识别流程进行整体测试。

该测试用于验证系统在连续短句场景下的端到端识别能力，
包括样本拼接、滑动窗口识别、词段合并、候选序列生成以及受约束语义修正等环节。

本节采用 core61 测试集作为拼接句子视频测试样例。

该测试集共包含 61 条核心短句，
短句长度覆盖 2 至 6 个 Gloss 不等，
能够反映系统在不同长度连续动作序列下的识别表现。

需要说明的是，
由于公开可用的单词级手语视频资源有限，
当前实验采用的是 25 类小词表模型。

词表规模和单词样本数量会直接影响句子级识别效果。

因此，
本节的重点并不是将拼接句子识别结果理解为大规模通用手语识别能力，
而是验证在当前小词表和有限样本条件下，
系统是否能够跑通从连续视频到 Gloss 序列输出的完整链路，
以及候选序列生成和语义修正是否具有进一步提升空间。

为了避免单次随机划分或单次运行结果对实验结论造成偏倚，
本文不以某一个随机种子下的结果作为主结论，
而是统计 10 组随机种子的平均结果。

对于每条测试样例，
系统首先基于词级分类模型输出原始识别序列 \texttt{rawSequence}。

随后，
系统根据每个词段的 Top-K 候选结果生成候选 Gloss 序列集合。

最后，
DeepSeek 在候选集合中选择一个候选编号，
后端再根据该候选编号得到修正后的 Gloss 序列。

表~\ref{tab:chapter06-spliced-sentence-multi-seed-result} 给出了 10 组随机种子下的平均测试结果。

\begin{longtable}{L{0.32\textwidth} L{0.22\textwidth} L{0.36\textwidth}}
  \caption{拼接句子视频多随机种子平均结果表}
  \label{tab:chapter06-spliced-sentence-multi-seed-result}\\
  \toprule
  指标 & 平均结果 & 说明 \\
  \midrule
  total & 61 & 每组随机种子下的测试样例总数 \\
  raw\_exact & 35.3 / 61 & 原始检测序列平均完全匹配数量 \\
  raw\_accuracy & 57.87\% & 原始检测序列平均完全匹配率 \\
  candidate\_oracle\_exact & 47.1 / 61 & 候选集合中平均包含正确序列的样例数 \\
  candidate\_oracle\_accuracy & 77.21\% & 候选集合平均理论完全匹配率 \\
  selected\_exact & 36.8 / 61 & DeepSeek 候选重排序后平均完全匹配数量 \\
  selected\_accuracy & 60.33\% & DeepSeek 候选重排序后平均完全匹配率 \\
  net\_gain & +1.5 & DeepSeek 候选重排序带来的平均净提升样例数 \\
  \bottomrule
\end{longtable}

由表~\ref{tab:chapter06-spliced-sentence-multi-seed-result} 可知，
在 10 组随机种子的平均结果中，
原始拼接句子视频识别的完全匹配率为 57.87\%。

这说明在当前小词表条件下，
系统已经能够在一部分连续短句样例中生成完全正确的 Gloss 序列。

同时，
该结果也说明原始识别仍然受到训练资源规模、动作边界、局部误判和词段后处理策略的影响。

DeepSeek 候选重排序后的平均完全匹配率为 60.33\%，
相比原始识别结果有一定提升。

候选集合理论完全匹配率为 77.21\%，
明显高于原始识别结果和 DeepSeek 实际选择结果。

这说明，
虽然原始 Top-1 序列中存在一定错误，
但系统通过保留词段 Top-K 候选并组合生成候选序列，
能够在相当一部分错误样例中保留正确答案。

因此，
候选序列生成模块本身是有效的，
它证明了当前识别链路中存在较大的可恢复空间。

在这个意义上，
候选集合理论上限并不是最终系统效果，
而是用于说明：
只要候选排序策略进一步优化，
句子级识别结果仍有继续提升的可能。

图~\ref{fig:chapter06-core61-average-summary} 进一步展示了 core61 句子集在多随机种子下的平均完全匹配准确率。

\WhuAutoFigure{0.82\textwidth}{0.42\textheight}{figures/chapter06/core61-average-summary.png}{Core61 句子集平均识别准确率汇总}{fig:chapter06-core61-average-summary}

由图~\ref{fig:chapter06-core61-average-summary} 可以看出，
候选上限明显高于原始识别和 DeepSeek 候选重排序结果。

这表明当前候选序列生成方法能够有效扩大正确答案空间。

同时，
DeepSeek 候选重排序结果高于原始识别结果，
说明受约束语义修正策略具有一定的稳定增益。

但是，
DeepSeek 实际选择结果与候选上限之间仍存在差距。

该差距并不意味着语义修正机制无效，
而是说明在当前小词表、小样本和候选信息有限的条件下，
候选排序仍然是系统后续优化的重要方向。

图~\ref{fig:chapter06-core61-accuracy-comparison} 展示了不同随机种子下原始识别、DeepSeek 候选重排序和候选上限三类结果的对比。

\WhuAutoFigure{0.92\textwidth}{0.48\textheight}{figures/chapter06/core61-accuracy-comparison.png}{Core61 句子集多随机种子识别准确率对比}{fig:chapter06-core61-accuracy-comparison}

从图~\ref{fig:chapter06-core61-accuracy-comparison} 可以看出，
不同随机种子下三类结果的整体趋势基本一致。

候选上限始终明显高于原始识别结果，
说明 Top-K 候选保留和候选序列生成策略具有稳定价值。

DeepSeek 候选重排序结果整体略高于原始识别结果，
说明受约束语义修正能够在不同随机种子下保持一定正向作用。

本文前期也尝试过更激进的语义修正策略。

但是，
在当前 25 类小词表和有限训练样本条件下，
激进策略容易过度依赖语言合理性，
从而产生脱离视觉证据的改写结果。

这种情况下，
修正后的序列反而可能低于原始识别序列。

因此，
本文最终采用候选集合约束下的保守重排序策略。

该策略的目标不是让大语言模型自由改写识别结果，
而是在尽量保持视觉证据忠实性的前提下，
从有限候选中选择更合理的序列。

从连续 10 组随机种子的结果看，
该策略能够带来稳定但有限的平均提升，
说明 DeepSeek 语义修正在当前系统中具有一定潜力。

结合实验结果可以归纳出，
拼接句子视频识别中主要存在以下几类错误现象。

表~\ref{tab:chapter06-sentence-error-types} 给出了句子视频识别错误现象归纳。

\begin{longtable}{L{0.22\textwidth} L{0.34\textwidth} L{0.36\textwidth}}
  \caption{句子视频识别错误现象归纳表}
  \label{tab:chapter06-sentence-error-types}\\
  \toprule
  错误类型 & 典型表现 & 可能原因 \\
  \midrule
  插入误报 & 序列中多出无关 Gloss，如两个动作之间出现额外词段 & 动作边界或过渡帧被误识别为具体词汇 \\
  漏词 & 期望序列中的某个 Gloss 未出现在检测序列中 & 词段置信度不足、动作持续时间短或被 NMS 抑制 \\
  重复或冗余输出 & 同一语义片段附近出现重复词或冗余词 & 滑动窗口重叠导致同一动作被多个窗口重复响应 \\
  局部替换错误 & 某个 Gloss 被识别为动作相近或语义无关的其他词 & 训练样本不足、手型相近或动作轨迹相似，导致词级分类混淆 \\
  候选排序未命中 & 正确序列存在于候选集合中，但最终未被 DeepSeek 选中 & 候选排序阶段对视觉证据、置信度信息和语义合理性的综合判断仍不充分 \\
  \bottomrule
\end{longtable}

由表~\ref{tab:chapter06-sentence-error-types} 可知，
插入误报和冗余输出是拼接句子视频识别中的重要问题。

这与滑动窗口识别机制的特点有关。

为了覆盖连续动作中的词汇边界，
系统需要使用相互重叠的窗口进行密集预测。

这种方式提高了动作片段被捕获的概率，
但也会使词与词之间的过渡帧形成额外候选窗口，
从而引入边界插入词或冗余输出。

漏词和局部替换错误则更多反映了词级分类模型和后处理策略的局限。

由于高质量单词级视频资源较少，
当前词级模型的训练样本规模有限。

当某个词段持续时间较短、
关键点检测质量不足，
或对应类别与其他动作类别相似时，
该词段可能被过滤、被抑制，
或被识别为其他 Gloss。

候选排序未命中则说明，
候选集合中存在正确答案并不等于最终系统一定能够选中正确答案。

在当前策略下，
DeepSeek 被限制在程序生成的候选集合中进行选择，
这能够避免自由生成带来的不受控改写，
但也要求候选信息本身能够充分表达视觉证据和局部置信度。

如果候选描述不足，
或原始序列在自然语言层面仍然可以解释，
DeepSeek 可能倾向于保留较保守的候选。

总体来看，
拼接句子视频整体测试表明，
当前系统已经能够完成从连续视频到 Gloss 序列的端到端识别。

候选集合的理论命中率显著高于原始识别完全匹配率，
说明 Top-K 候选保留和候选序列生成具有继续优化价值。

DeepSeek 候选重排序带来了稳定但有限的平均提升，
说明受约束语义修正在当前小词表条件下具有一定潜力。

同时，
实验结果也说明当前系统的主要限制来自两个方面。

一方面，
高质量单词级手语视频资源不足，
导致词表规模较小，
词级模型对相似动作的区分能力仍有提升空间。

另一方面，
候选排序阶段尚未充分利用候选集合中已经存在的正确序列。

后续若要进一步提升句子级识别效果，
需要继续扩充高质量训练资源，
提高词级模型的分类能力，
并在保持候选约束的前提下改进语义修正模块对视觉证据和语言合理性的综合判断。
% !TEX root = ../../../main.tex

\subsection{典型识别样例分析}


\WhuAutoFigure{0.82\textwidth}{0.42\textheight}{figures/chapter06/typical-sentence-result.png}{典型句子视频识别结果展示图}{fig:chapter06-typical-sentence-result}


\begin{longtable}{L{0.15\textwidth} L{0.17\textwidth} L{0.17\textwidth} L{0.15\textwidth} L{0.17\textwidth} L{0.14\textwidth}}
  \caption{典型识别样例分析表}
  \label{tab:chapter06-typical-case-analysis}\\
  \toprule
  样例 & 期望 Gloss & 识别 Gloss & 中文直译 & 语义修正结果 & 分析说明 \\
  \midrule
  正确识别样例 & 待补充 & 待补充 & 待补充 & 待补充 & 待补充 \\
  插入误报样例 & 待补充 & 待补充 & 待补充 & 待补充 & 待补充 \\
  局部替换样例 & 待补充 & 待补充 & 待补充 & 待补充 & 待补充 \\
  \bottomrule
\end{longtable}

% !TEX root = ../../../main.tex

\subsection{DeepSeek 语义修正效果分析}

在拼接句子视频识别流程中，
DeepSeek 并不直接处理视频帧，
而是在后端生成的候选序列集合中选择一个候选编号。

因此，
本节重点分析 DeepSeek 在受约束候选空间中的重排序效果，
而不是将其视为自由生成式翻译模块。

根据 10 组随机种子的平均结果，
原始识别序列的平均完全匹配率为 57.87\%，
DeepSeek 候选重排序后的平均完全匹配率为 60.33\%，
候选集合理论上限为 77.21\%。

这说明 DeepSeek 候选重排序能够带来一定提升，
但实际选择结果与候选集合理论上限之间仍然存在明显差距。

表~\ref{tab:chapter06-deepseek-effect-summary} 给出了 DeepSeek 候选重排序效果汇总。

\begin{longtable}{L{0.28\textwidth} L{0.22\textwidth} L{0.40\textwidth}}
  \caption{DeepSeek 候选重排序效果汇总表}
  \label{tab:chapter06-deepseek-effect-summary}\\
  \toprule
  指标 & 结果 & 说明 \\
  \midrule
  raw\_accuracy & 57.87\% & 不经过 DeepSeek 重排序时，原始识别序列的平均完全匹配率 \\
  selected\_accuracy & 60.33\% & DeepSeek 从候选集合中选择后得到的平均完全匹配率 \\
  candidate\_oracle\_accuracy & 77.21\% & 候选集合中存在正确序列时的理论上限 \\
  accuracy\_gain & +2.46 个百分点 & DeepSeek 候选重排序相对原始识别的平均准确率提升 \\
  net\_gain & +1.50 条 & DeepSeek 候选重排序平均多修正的样例数量 \\
  oracle\_gap & 16.88 个百分点 & DeepSeek 实际选择结果与候选上限之间的差距 \\
  \bottomrule
\end{longtable}

由表~\ref{tab:chapter06-deepseek-effect-summary} 可知，
DeepSeek 候选重排序相对原始识别结果带来了平均 2.46 个百分点的准确率提升。

从样例数量角度看，
该模块平均能够带来 1.50 条样例的净提升。

虽然提升幅度并不大，
但该结果是在候选集合约束和视觉证据约束下取得的。

也就是说，
DeepSeek 并没有被允许自由新增词汇、重排词序或生成候选集合之外的 Gloss，
而是只能在后端生成的有限候选序列中进行选择。

因此，
该结果更适合被理解为一种稳定但有限的受约束修正收益。

图~\ref{fig:chapter06-core61-deepseek-net-gain} 展示了 DeepSeek 候选重排序在 10 组随机种子下带来的净收益。

\WhuAutoFigure{0.86\textwidth}{0.44\textheight}{figures/chapter06/core61-deepseek-net-gain.png}{DeepSeek 候选重排序净收益统计}{fig:chapter06-core61-deepseek-net-gain}

由图~\ref{fig:chapter06-core61-deepseek-net-gain} 可以看出，
不同随机种子下 DeepSeek 候选重排序带来的净收益存在一定波动，
但整体平均净收益为正。

这说明当前语义修正策略并非只在单次实验中偶然有效，
而是在多组随机种子下都表现出一定的稳定增益。

需要说明的是，
本文采用的 DeepSeek 语义修正策略并非一开始即设定为当前形式，
而是在多轮实验对比后逐步收敛得到的结果。

在前期实验中，
系统曾尝试让大语言模型进行更激进的语义纠错，
例如允许其根据自然语言合理性对原始 Gloss 序列进行较大幅度改写。

但是，
在当前小词表和小样本条件下，
这类激进策略容易产生过度修正问题。

一方面，
当前词级识别模型支持的 Gloss 类别数量有限，
词表规模较小，
大语言模型的自然语言先验可能强于视觉证据约束，
从而倾向于生成语义更自然但不一定受到视频识别结果支持的表达。

另一方面，
单个词语视频资源数量较少，
训练样本覆盖不足，
模型在相似动作之间仍存在混淆。

如果语义修正阶段过度依赖语言合理性，
反而可能放大局部误判，
导致修正后结果低于原始识别结果。

因此，
本文最终采用受候选集合约束的保守修正策略。

在该策略下，
DeepSeek 只能从后端生成的候选序列中选择一个候选编号，
不能自由新增词汇，
不能重排词序，
也不能生成候选集合之外的 Gloss。

这种方式虽然限制了语义修正幅度，
但能够保证最终结果始终具有视觉识别证据来源，
避免大语言模型自由改写导致结果失真。

从连续 10 组随机种子的实验结果看，
该策略相较原始识别结果具有稳定但有限的平均提升，
说明 DeepSeek 语义修正在当前系统中具有一定潜力。

与此同时，
候选集合理论上限明显高于 DeepSeek 实际选择结果，
说明候选序列生成模块已经在一定程度上保留了正确答案空间，
但候选排序阶段仍未充分利用这些候选。

表~\ref{tab:chapter06-deepseek-correction-phenomena} 给出了 DeepSeek 语义修正中的主要现象归纳。

\begin{longtable}{L{0.24\textwidth} L{0.34\textwidth} L{0.34\textwidth}}
  \caption{DeepSeek 语义修正现象归纳表}
  \label{tab:chapter06-deepseek-correction-phenomena}\\
  \toprule
  现象类型 & 典型表现 & 说明 \\
  \midrule
  保留原始序列 & 原始识别序列已经较为连贯时，DeepSeek 倾向于选择原始候选 & 有助于避免破坏正确样例，但也可能导致部分插入误报未被删除 \\
  删除插入词 & 原始序列中存在明显边界噪声词时，DeepSeek 可能选择删除候选 & 能够修正部分插入误报，是当前语义修正的主要收益来源之一 \\
  局部替换 & 当 Top-K 候选中存在更合适的词段标签时，DeepSeek 可能选择替换候选 & 可修正部分局部分类错误，但替换次数受候选生成规则限制 \\
  无法恢复漏词 & 原始识别或候选集合中缺少某个关键 Gloss 时，最终结果仍然缺词 & 体现了候选约束机制的边界，DeepSeek 不会凭空新增无视觉证据的词 \\
  候选存在但未选中 & 正确序列在候选集合中，但 DeepSeek 最终选择其他候选 & 说明候选排序仍是当前系统的主要优化方向 \\
  \bottomrule
\end{longtable}

由表~\ref{tab:chapter06-deepseek-correction-phenomena} 可知，
DeepSeek 语义修正模块在当前系统中主要发挥受约束候选重排序作用。

它可以删除部分明显的插入词，
也可以在少数情况下完成局部替换。

但是，
它并不能替代视觉识别模块，
也不能突破候选集合本身的限制。

如果正确 Gloss 没有出现在词段 Top-K 候选或候选序列集合中，
DeepSeek 不会凭空生成该词。

这种设计在一定程度上牺牲了自由修正能力，
但换来了更强的视觉忠实性和结果稳定性。

从自然语言表达角度看，
DeepSeek 还承担了将 Gloss 序列转换为用户可读中文的任务。

对于语义关系明确的短句，
DeepSeek 可以将电报式 Gloss 序列转换为较自然的中文句子。

例如，
\texttt{what you want today} 可以转换为“你今天想要什么？”。

对于语义关系不完整的短句，
系统会通过低置信度或不完整提示保留不确定性，
而不是强行补全为完整句子。

这种方式能够降低自然化翻译阶段的过度补全风险。

综合来看，
DeepSeek 语义修正模块在当前系统中具有两个主要价值。

第一，
它在候选集合约束下能够带来稳定但有限的句子级识别增益。

第二，
它能够将 Gloss 序列转换为更符合用户阅读习惯的中文表达，
提升系统输出的可理解性。

同时，
实验结果也表明，
当前系统的主要瓶颈并不只是语义修正策略本身。

由于公开可用的单词级手语视频资源有限，
当前模型只能在较小词表上进行训练和测试。

词表规模较小会限制句子级组合表达能力，
样本数量不足也会影响词级模型对相似动作的区分能力。

在这种条件下，
过于激进的语义修正策略容易脱离视觉证据，
而受候选约束的保守策略更适合作为当前阶段的系统方案。

后续优化应重点从三个方面展开。

一是继续扩充高质量单词级视频资源，
提高训练样本数量和动作覆盖范围。

二是提升词级分类模型对相似手型、姿态和运动轨迹的区分能力，
减少原始识别阶段的局部误判和漏词。

三是在保持候选约束的前提下，
改进候选排序提示词、候选特征描述和视觉证据表达方式，
使 DeepSeek 能够更充分地利用候选集合中已经存在的正确序列。

这样既能保持识别结果对视觉证据的忠实性，
又能进一步发挥语义修正模块在句子级识别中的潜力。

% !TEX root = ../../main.tex

\section{连续识别方案探索与局限性}

除采用“词级分类模型 + 滑动窗口 + 词段后处理”方案外，本项目在实验过程中也对更直接的连续手语识别方案进行了探索。
连续识别方案的基本思想是让模型直接从一段较长的视频特征序列中输出 Gloss 序列，而不是先将视频切分为大量候选窗口，再通过分类和后处理得到词段。
这类方案通常需要使用 CTC 等序列建模方法，以处理输入帧序列和输出 Gloss 序列长度不一致的问题。

在探索阶段，项目曾基于 CE-CSL 数据进行 CTC 方向实验。
一类实验采用 Transformer Encoder + CTC 结构，输入特征使用 raw\_delta 模式，并通过位置编码、padding mask 和 CTC 解码完成序列预测。
该方案能够完成训练并输出序列，但实验结果并不理想，主要问题表现为预测序列偏短、删除错误较多，说明模型虽然具备一定学习能力，但在当前特征、数据规模和训练配置下，仍难以稳定输出完整 Gloss 序列。
因此，该方向没有被纳入 HearBridge 当前系统主线。

另一类探索使用 Conformer Encoder + CTC 结构，并在 CE-CSL 数据上引入 controlled vocabulary，将低频 Gloss 映射为 \texttt{<unk>}，以降低完整词表带来的稀疏性问题。
该方案在训练脚本中记录 train\_loss、dev\_loss 和 dev\_TER，并保存最佳 dev\_TER、最佳 dev\_loss 和最后一轮 checkpoint，适合作为后续连续识别研究的实验基础。
但由于该实验依赖较大规模连续手语数据、训练成本较高，且与 HearBridge 当前 25 词短句识别任务和系统接入链路并不完全一致，因此本项目最终没有将其作为毕业设计系统的主要实现方案。

相比之下，滑动窗口词级识别方案更适合本项目阶段目标。
一方面，WLASL 小词表能够较方便地构建 $20\times232$ 词级训练样本，模型训练和调试成本较低；
另一方面，滑动窗口方案可以在没有精确句子级边界标注的情况下，对短视频中的局部词段进行密集识别，并通过 NMS 和 DeepSeek 候选排序降低插入误报和语义不自然问题。

当前方案仍然存在明显局限。
首先，滑动窗口方式依赖固定窗口长度和步长，动作边界仍然通过后处理间接确定，无法像完整连续识别模型那样从序列层面统一建模。
其次，词级分类模型的识别范围受限于小词表，无法覆盖开放场景下的大规模 Gloss。
再次，DeepSeek 语义修正只能在程序生成的候选集合内选择更合理的序列，不能弥补视觉模型完全未识别出的关键词。

% !TEX root = ../../main.tex

\section{本章小结}

本章围绕 HearBridge 系统中的手势识别服务实现、句子视频识别流程与实验结果分析展开说明。

首先，
针对移动端实时训练与识别场景，
本章介绍了自采样本采集链路、手势特征构造、词级模型训练、模型版本发布和实时识别闭环。

在该流程中，
移动端负责采集用户手势图像帧并通过 WebSocket 发送至 Python 识别服务，
Python 端完成 MediaPipe 关键点检测、特征提取、样本保存、模型训练和实时预测，
后台管理端则负责样本管理、模型训练触发、模型版本发布和发布模型重载。

该部分验证了系统能够完成从移动端样本采集、服务端特征转换、模型训练发布到移动端实时识别展示的完整闭环。

其次，
针对句子视频识别场景，
本章说明了基于 WLASL 小词表数据的词级分类模型构建过程。

系统采用 $20\times232$ 的 plus 特征作为词级分类输入，
并在 25 类小词表上完成模型训练与测试。

实验结果表明，
词级模型虽然在 Top-1 分类准确率上仍有提升空间，
但 Top-3 准确率较高，
能够为后续词段 Top-K 候选保留、候选序列生成和语义修正提供有效基础。

在句子视频识别流程方面，
本章进一步介绍了上传视频后的帧级检测缓存、滑动窗口输入生成、逐窗口词级分类、词段合并、NMS 抑制、Gloss 序列生成和 DeepSeek 语义修正机制。

该流程将长度可变的短句视频转换为多个固定长度候选窗口，
再通过词级模型得到窗口预测结果，
最终经词段后处理生成原始 Gloss 序列。

随后，
后端基于词段 Top-K 候选生成有限候选序列，
并约束 DeepSeek 只能从候选集合中选择一个候选编号，
从而在保持视觉证据约束的前提下完成语义修正和中文自然化表达。

在实验分析部分，
本章对 core61 拼接句子视频测试集进行了多随机种子实验。

结果表明，
原始句子级识别已经能够完成部分短句的完整匹配；
候选集合理论上限明显高于原始识别结果，
说明 Top-K 候选保留和候选序列生成具有较大的可恢复空间；
DeepSeek 候选重排序相比原始结果具有稳定但有限的平均提升，
说明受约束语义修正在当前小词表条件下具有一定潜力。

同时，
实验也表明当前系统仍存在插入误报、漏词、冗余输出、局部替换错误和候选排序未命中等问题。

这些问题一方面来自滑动窗口机制对动作边界的间接建模，
另一方面也与单词级训练资源不足、词表规模较小和相似动作类别混淆有关。

因此，
当前方案更适合作为有限词表短句识别和系统功能验证方案，
尚不能等同于开放域连续手语识别模型。

最后，
本章还对更直接的连续手语识别方案进行了探索与说明。

相比词级分类与滑动窗口方案，
基于 Transformer、Conformer 和 CTC 的连续识别方法在任务形式上更接近真实连续手语识别，
能够直接从长视频特征序列中预测 Gloss 序列。

但是，
该方向对连续数据质量、序列建模能力、训练策略、参数调优和解码方法要求更高。

在当前毕业设计周期和系统交付目标下，
相关实验尚未达到可稳定接入系统的程度。

因此，
本项目最终采用训练难度更可控、工程链路更稳定的词级分类与滑动窗口方案，
优先完成可运行、可测试、可展示的端到端识别系统。

综上，
本章完成了 HearBridge 手势识别服务从移动端实时识别到句子视频识别的核心实现与实验验证。

当前系统已经能够支撑样本采集、模型训练、模型发布、实时识别、句子视频上传识别和中文语义表达等主要功能。

后续工作可从扩充高质量手语视频资源、扩大词表规模、优化词级模型特征表达、改进词段边界处理、增强候选排序能力以及继续探索连续识别模型等方面进一步提升系统性能。
